% ROZDZIAL 5 -- Concepts i requires w C++20
% Fragment do wlaczenia do glownego dokumentu po rozdziale 4.
% Wymagane w preambule: pakiet listings i ustawienie language=C++.
% Uzycie: % ROZDZIAL 5 -- Concepts i requires w C++20
% Fragment do wlaczenia do glownego dokumentu po rozdziale 4.
% Wymagane w preambule: pakiet listings i ustawienie language=C++.
% Uzycie: % ROZDZIAL 5 -- Concepts i requires w C++20
% Fragment do wlaczenia do glownego dokumentu po rozdziale 4.
% Wymagane w preambule: pakiet listings i ustawienie language=C++.
% Uzycie: % ROZDZIAL 5 -- Concepts i requires w C++20
% Fragment do wlaczenia do glownego dokumentu po rozdziale 4.
% Wymagane w preambule: pakiet listings i ustawienie language=C++.
% Uzycie: \input{rozdzial_5_concepts_cpp.tex}

\chapter[Concepts i requires w C++]{Koncepcje i ograniczenia szablonów (\textit{concepts} i \texttt{requires})}
\label{chap:concepts-cpp}

\section{Po co wprowadzono koncepcje?}

Szablony C++ pozwalają pisać kod niezależny od konkretnego typu.
Tradycyjny szablon deklaruje jednak zwykle \emph{co kompilator ma zrobić},
nie wyjaśniając wprost, \emph{jakich własności} wymaga od argumentów.
Jeśli operacja jest niepoprawna, błąd może zostać wykryty dopiero w ciele
szablonu, a komunikat diagnostyczny bywa długi.

\begin{lstlisting}
template<class T>
T twice(const T& x) {
    return x + x;           // requires operator+
}

// twice(7);                // OK
// twice(SomeType{});       // may fail inside function body
\end{lstlisting}

Od C++20 można umieścić wymagania na poziomie interfejsu:

\begin{lstlisting}
#include <concepts>

template<class T>
concept Addable = requires(const T& a, const T& b) {
    { a + b } -> std::convertible_to<T>;
};

template<Addable T>
T twice(const T& x) {
    return x + x;
}
\end{lstlisting}

Koncepcja (\textit{concept}) to nazwany predykat sprawdzany podczas
kompilacji, opisujący wymagania wobec argumentów szablonu.
Może dotyczyć istnienia operatorów, typów składowych, własności typów,
kategorii iteratorów i wielu innych cech. \emph{Constraints}
(ograniczenia) wpływają zarówno na dopuszczalność konkretyzacji,
jak i na wybór przeciążenia funkcji.

\paragraph{Wersja języka.}
W całym rozdziale zakładamy \textbf{C++20} oraz kompilację
np. poleceniem \texttt{g++ -std=c++20 -Wall -Wextra -Wpedantic main.cpp}.
Definicje koncepcji standardowych znajdują się przede wszystkim
w nagłówkach \texttt{<concepts>}, \texttt{<iterator>} i
\texttt{<ranges>}.

\section{Pierwsza koncepcja: składnia i zastosowanie}
\label{sec:concepts-pierwsza}

Definicja nazwanej koncepcji ma postać
\texttt{template<...> concept Nazwa = warunek;}. Warunek powinien
być wyrażeniem stałym typu \texttt{bool}. Definicja koncepcji
występuje na poziomie przestrzeni nazw; nie jest zwykłą funkcją.

\begin{lstlisting}
#include <concepts>
#include <type_traits>

template<class T>
concept Integer = std::is_integral_v<T>;

static_assert(Integer<int>);
static_assert(Integer<bool>);    // bool is integral
static_assert(!Integer<double>);
\end{lstlisting}

Koncepcję można sprawdzić wyrażeniem \texttt{Integer<int>} lub
\texttt{static\_assert}, ale nie należy jej mylić ze zwykłą
zmienną. W szczególności \texttt{concept} nie daje się jawnie
specjalizować w sposób analogiczny do szablonu klasy.

\subsection{Cztery podstawowe sposoby ograniczania funkcji}

Poniższe deklaracje opisują ten sam warunek: funkcja jest przeznaczona
dla typów całkowitych w sensie \texttt{std::integral}.

\begin{lstlisting}
#include <concepts>

template<class T>
requires std::integral<T>
T square1(T x) { return x * x; }

template<class T>
T square2(T x) requires std::integral<T> {
    return x * x;
}

template<std::integral T>
T square3(T x) { return x * x; }

std::integral auto square4(std::integral auto x) {
    return x * x;
}
\end{lstlisting}

W ostatnim przykładzie ograniczony jest nie tylko typ argumentu,
ale też \emph{typ wyniku}. Sama dedukcja typu wyniku nadal odbywa
się na podstawie wyrażenia \texttt{x * x}; koncept jest sprawdzany
po jej przeprowadzeniu. Zapis \texttt{std::integral auto} w parametrze
jest skrótem szablonu funkcji, tzw. \emph{abbreviated function template}.

\paragraph{Kiedy którą postać wybrać?}
Dla jednego prostego ograniczenia zapis z ograniczonym
parametrem szablonu, np. \texttt{template<\allowbreak std::integral T>},
jest krótki i czytelny. Jeśli zależność dotyczy kilku parametrów,
najwygodniejsza bywa klauzula \texttt{requires}.
Skrócona postać z \texttt{auto} przydaje się w niewielkich funkcjach
i lambdach.

\section{Dwie role słowa \texttt{requires}}

\begin{center}\small
\begin{tabularx}{\textwidth}{@{}p{0.34\textwidth}X@{}}
\toprule
Zapis & Za co odpowiada \\
\midrule
\code{concept C = ...} & Nadaje nazwę logicznemu ograniczeniu parametrów szablonu. \\
\code{requires C<T>} & Ogranicza udział szablonu w przeciążaniu do typów spełniających \code{C}. \\
\code{requires(T x) { x.f(); }} & Sprawdza poprawność wskazanego wyrażenia dla podstawionego typu. \\
\code{requires requires(...) { ... }} & Pierwsze \code{requires} wprowadza ograniczenie, drugie wyrażenie sprawdzające. \\
\code{static_assert(C<T>)} & Wymusza spełnienie warunku w danym miejscu i zgłasza błąd kompilacji, gdy jest fałszywy. \\
\bottomrule
\end{tabularx}
\end{center}
\label{sec:concepts-dwa-requires}

Identyczny zapis słowny \texttt{requires} występuje w dwóch
różnych konstrukcjach:
\begin{enumerate}
  \item \emph{requires-clause} -- klauzula ograniczająca deklarację;
  \item \emph{requires-expression} -- wyrażenie logiczne, które bada
  poprawność podanych wymagań.
\end{enumerate}

\begin{lstlisting}
#include <concepts>

template<class T>
concept Incrementable = requires(T value) { // requires-expression
    ++value;
};

template<class T>
requires Incrementable<T>                // requires-clause
void increment(T& x) {
    ++x;
}
\end{lstlisting}

Można je połączyć bez nadawania koncepcji osobnej nazwy:

\begin{lstlisting}
template<class T>
requires requires(T x) { ++x; }
void increment_again(T& x) {
    ++x;
}
\end{lstlisting}

Pierwsze \texttt{requires} w ostatnim przykładzie zaczyna klauzulę;
drugie wprowadza wyrażenie sprawdzające wymagania. To poprawna
składnia, choć przy wielokrotnym użyciu lepiej utworzyć nazwaną
koncepcję.

\section{Cztery rodzaje wymagań w \texttt{requires-expression}}

\begin{center}\small
\begin{tabularx}{\textwidth}{@{}p{0.31\textwidth}X@{}}
\toprule
Rodzaj wymagania & Co sprawdza \\
\midrule
Proste: \code{a + b;} & Czy wyrażenie jest poprawne składniowo dla danych typów. \\
Typowe: \code{typename T::value_type;} & Czy wskazany typ zależny istnieje. \\
Złożone: \code{{ f(x) } -> C;} & Czy wyrażenie jest poprawne oraz czy jego typ spełnia ograniczenie \code{C}. \\
Zagnieżdżone: \code{requires C<T>;} & Czy dodatkowy warunek logiczny jest spełniony. \\
\bottomrule
\end{tabularx}
\end{center}
\label{sec:concepts-cztery-wymagania}

Wyrażenie \texttt{requires} umożliwia opisanie czterech rodzajów
warunków. Parametry wpisane w nawiasach \texttt{requires(T x)}
są jedynie \emph{lokalnymi parametrami do sprawdzania wyrażeń};
nie tworzy się dla nich obiektów podczas działania programu.

\subsection{Wymaganie proste (\textit{simple requirement})}

\begin{lstlisting}
#include <ostream>

template<class T>
concept Printable = requires(std::ostream& os, const T& x) {
    os << x;                    // must be well-formed
};
\end{lstlisting}

Wyrażenie \texttt{os << x}
nie jest tutaj wykonywane; kompilator sprawdza jego poprawność.
Proste wymaganie nie nakłada warunku na typ zwracany.

\subsection{Wymaganie typu (\textit{type requirement})}

\begin{lstlisting}
#include <vector>

template<class T>
concept HasValueType = requires {
    typename T::value_type;    // nested type must exist
};

static_assert(HasValueType<std::vector<int>>);
static_assert(!HasValueType<int>);
\end{lstlisting}

Słowo \texttt{typename} jest tu składnikiem składni wymagania,
a nie deklaracją nowego typu.

\subsection{Wymaganie złożone (\textit{compound requirement})}

\begin{lstlisting}
#include <concepts>

template<class T>
concept AddableNoThrow = requires(const T& a, const T& b) {
    { a + b } noexcept -> std::same_as<T>;
};
\end{lstlisting}

Zapis \texttt{\{ wyrazenie \} noexcept -> Koncepcja;}
może sprawdzać jednocześnie: poprawność wyrażenia,
to, że jest ono \texttt{noexcept}, i własność typu
\texttt{decltype((wyrazenie))}. Obie części, \texttt{noexcept}
oraz \texttt{-> Koncepcja}, są opcjonalne.

\begin{lstlisting}
template<class T>
concept PlusReturnsT = requires(const T& a, const T& b) {
    { a + b } -> std::same_as<T>;
};

static_assert(PlusReturnsT<int>);
// A type with a proxy return value may not satisfy this.
\end{lstlisting}

\paragraph{Uwaga na typ i konwersje.}
Koncepcja \texttt{std::same\_as<T>} żąda \emph{dokładnie}
typu \texttt{T}, natomiast koncepcja
\texttt{std::convertible\_\allowbreak to<T>} dopuszcza konwersję.
Dla \texttt{int} operator \texttt{+} daje \texttt{int}, ale np.
dla małych typów całkowitych może zajść promocja do \texttt{int}.

\subsection{Wymaganie zagnieżdżone (\textit{nested requirement})}

\begin{lstlisting}
#include <type_traits>

template<class T>
concept SmallTrivial = requires {
    requires sizeof(T) <= 16;
    requires std::is_trivially_copyable_v<T>;
};

static_assert(SmallTrivial<int>);
\end{lstlisting}

Wymaganie zagnieżdżone zaczyna się kolejnym \texttt{requires}
i musi zawierać warunek stały. Jest wygodne, gdy wymaganie
dotyczy własności typu, a nie określonego wyrażenia.

\subsection{Przykład: połączenie różnych wymagań}

\begin{lstlisting}
#include <concepts>
#include <cstddef>

template<class T>
concept Indexable = requires(T& obj, std::size_t i) {
    typename T::value_type;
    { obj.size() } -> std::convertible_to<std::size_t>;
    obj[i];
};
\end{lstlisting}

To przykład \emph{minimalnego kontraktu składniowego}.
Koncepcja nie dowodzi, że indeks \texttt{i} jest w zakresie,
że operacja jest szybka ani że wynik indeksowania zachowuje się
jak zwykła referencja. Takie własności wymagają dodatkowych założeń.

\section{Standardowe koncepcje z biblioteki C++20}
\label{sec:concepts-standardowe}

Przed napisaniem własnej koncepcji warto sprawdzić, czy istnieje
gotowa w standardowej bibliotece.

\subsection{Typy, konwersje i dziedziczenie}

\begin{lstlisting}
#include <concepts>

static_assert(std::same_as<int, int>);
static_assert(!std::same_as<int, long>);
static_assert(std::integral<unsigned>);
static_assert(std::signed_integral<int>);
static_assert(std::unsigned_integral<unsigned>);
static_assert(std::floating_point<double>);
static_assert(std::convertible_to<int, double>);

struct Base {};
struct Derived : Base {};
static_assert(std::derived_from<Derived, Base>);
\end{lstlisting}

\texttt{std::derived\_from<D,B>} obejmuje publiczne i jednoznaczne
dziedziczenie oraz przypadek tego samego typu. Różni się pod tym
względem od samego \texttt{std::is\_base\_of\_v}.

\subsection{Konstruowalność, kopiowanie, przenoszenie i porównywanie}

Przydatne są m.in. koncepcje wyrażające inicjalizowalność,
konstruowalność, możliwość przenoszenia i kopiowania oraz porównywania:
\begin{lstlisting}
default_initializable    constructible_from
move_constructible       copy_constructible
movable                  copyable
semiregular              regular
equality_comparable      totally_ordered
\end{lstlisting}

\begin{lstlisting}
#include <concepts>
#include <memory>
#include <string>

static_assert(std::regular<int>);
static_assert(std::copyable<std::string>);
static_assert(std::movable<std::unique_ptr<int>>);
static_assert(!std::copyable<std::unique_ptr<int>>);
static_assert(std::totally_ordered<int>);
\end{lstlisting}

Koncepcje te opisują interfejs i zakładają także określone
\emph{prawa semantyczne} (np. sensowne porównywanie). Sam kompilator
nie potrafi udowodnić dowolnej własności semantycznej implementacji.

\subsection{Funkcje wywoływalne i predykaty}

\begin{lstlisting}
#include <concepts>

struct IsPositive {
    bool operator()(int x) const { return x > 0; }
};

static_assert(std::invocable<IsPositive, int>);
static_assert(std::predicate<IsPositive, int>);
static_assert(!std::predicate<IsPositive, void*>);
\end{lstlisting}

\texttt{std::invocable<F,Args...>} sprawdza możliwość wywołania,
natomiast \texttt{std::predicate} wymaga dodatkowo wyniku,
którego można użyć jako warunku logicznego, i wiąże się
z określonymi wymaganiami semantycznymi.

\subsection{Iteratory i zakresy}

Koncepcje iteratorów są udostępniane przez \texttt{<iterator>},
a koncepcje zakresów przez \texttt{<ranges>}.

\begin{lstlisting}
#include <ranges>
#include <vector>
#include <list>

static_assert(std::ranges::random_access_range<
              std::vector<int>>);
static_assert(std::ranges::bidirectional_range<
              std::list<int>>);
static_assert(!std::ranges::random_access_range<
              std::list<int>>);
\end{lstlisting}

W szczególności przydatne są koncepcje \emph{input},
\emph{forward}, \emph{bidirectional}, \emph{random access}
i \emph{contiguous range}. Pozwalają opisywać wymagania
algorytmów bez uzależniania ich od konkretnego kontenera.

\section{Łączenie koncepcji i definiowanie własnych kontraktów}

Wymagania można łączyć operatorami logicznymi
\texttt{\&\&} i \texttt{||}. Często lepiej budować niewielkie,
nazwane koncepcje niż powtarzać jedno długie wyrażenie.

\begin{lstlisting}
#include <concepts>

template<class T>
concept Arithmetic = std::integral<T> || std::floating_point<T>;

template<class T>
concept NonBooleanInteger =
    std::integral<T> && !std::same_as<T, bool>;

template<Arithmetic T>
T cube(T x) { return x * x * x; }

static_assert(Arithmetic<double>);
static_assert(!NonBooleanInteger<bool>);
\end{lstlisting}

Dla połączenia \texttt{A \&\& B} wymaganie \texttt{B} jest
sprawdzane dopiero, gdy \texttt{A} jest spełnione; analogicznie
dla \texttt{A || B} drugi składnik jest potrzebny dopiero,
gdy pierwszy nie jest spełniony. Ta reguła chroni często przed
podstawianiem wyrażeń, które dla danego typu nie mają sensu.

\begin{lstlisting}
#include <concepts>
#include <type_traits>

template<class T>
concept HasIntegralValueType =
    requires { typename T::value_type; } &&
    std::integral<typename T::value_type>;
\end{lstlisting}

W powyższym przykładzie odwołanie do \texttt{T::value\_type}
po prawej stronie jest istotne dopiero po potwierdzeniu,
że typ składowy istnieje. Należy jednak pamiętać, że złożone
reguły podstawiania i normalizacji koncepcji mają swoje szczegóły;
nie każdą źle sformułowaną deklarację naprawi samo zwarcie logiczne.

\section{Koncepcje z wieloma parametrami}

Koncepcja nie musi opisywać własności jednego typu.
Może także opisywać \emph{relację między typami}.

\begin{lstlisting}
#include <concepts>

template<class From, class To>
concept SameType = std::same_as<From, To>;

template<class Derived, class Base>
concept PublicChildOf = std::derived_from<Derived, Base>;

template<class T, class U>
concept AddableTogether = requires(const T& a, const U& b) {
    a + b;
};

static_assert(AddableTogether<int, double>);
\end{lstlisting}

Wymienione koncepcje \texttt{SameType} i
\texttt{PublicChildOf} służą tu jedynie ilustracji składni.
W praktyce używaj bezpośrednio standardowych
\texttt{std::same\_as} i \texttt{std::derived\_from},
jeżeli nie dodajesz do nich istotnego znaczenia domenowego.

\paragraph{Wstawianie argumentu przy ograniczonym parametrze.}

\begin{lstlisting}
struct Animal {};

template<std::derived_from<Animal> T>
void inspect(const T&) {}

// Equivalent constraint: std::derived_from<T, Animal>
\end{lstlisting}

W zapisie powyżej typ \texttt{T} jest podstawiany
\emph{na pierwszej pozycji} koncepcji, zatem
chodzi o warunek z poniższego przykładu:
\begin{lstlisting}
std::derived_from<T, Animal>
\end{lstlisting}

\section{Przeciążanie funkcji i wybór bardziej szczegółowego ograniczenia}
\label{sec:concepts-przeciazanie}

Ograniczenia uczestniczą w wyborze przeciążenia. Jeśli dwie funkcje
są dopasowane do argumentów, kompilator może preferować tę,
której ograniczenie jest bardziej szczegółowe (ang.
\emph{subsumption}).

\begin{lstlisting}
#include <concepts>
#include <string_view>

template<class T>
requires std::integral<T>
constexpr std::string_view category(T) {
    return "integral";
}

template<class T>
requires std::signed_integral<T>
constexpr std::string_view category(T) {
    return "signed integral";
}

static_assert(category(12) == "signed integral");
static_assert(category(12u) == "integral");
\end{lstlisting}

\texttt{std::signed\_integral} jest zbudowane na
\texttt{std::integral} i dodaje dodatkowe wymaganie. Dzięki temu
bardziej szczegółowe przeciążenie zostaje wybrane dla \texttt{int}.

\subsection{Ważna pułapka: podobny zapis to nie to samo ograniczenie}

Porządek ograniczeń nie jest ogólnym dowodzeniem implikacji
matematycznych. Kompilator porównuje znormalizowane ograniczenia
atomowe i ich \emph{tożsamość}, a nie tylko sens pozornie podobnych
wyrażeń.

\begin{lstlisting}
#include <type_traits>

template<class T>
concept HasSize = requires(const T& x) { x.size(); };

template<class T>
concept HasSizeAndCopy =
    HasSize<T> && std::is_copy_constructible_v<T>;

template<class T> requires HasSize<T>
int classify(const T&) { return 1; }

template<class T> requires HasSizeAndCopy<T>
int classify(const T&) { return 2; }
\end{lstlisting}

Druga koncepcja \emph{wykorzystuje tę samą} koncepcję
\texttt{HasSize}, dlatego kompilator rozpoznaje wspólny
składnik. Jeśli powtórzymy ręcznie nawet identycznie wyglądające
wymaganie opisujące metodę \texttt{size()},
nie mamy automatycznej gwarancji identyczności ograniczeń
atomowych.

\paragraph{Rekomendacja.}
Twórz proste nazwane koncepcje i buduj bardziej szczegółowe
z ich wykorzystaniem. Ułatwia to konserwację kodu i
przewidywanie wyboru przeciążeń.

\section{Koncepcje dla szablonów klas i ich specjalizacji}

Ograniczać można nie tylko funkcje, lecz również szablony klas.

\begin{lstlisting}
#include <concepts>

template<std::integral T>
class Counter {
    T value_{};
public:
    void add(T delta) { value_ += delta; }
    T value() const { return value_; }
};

// Counter<int> a;       // OK
// Counter<double> b;    // ERROR: constraint not satisfied
\end{lstlisting}

Ograniczenia mogą też sterować wyborem częściowej specjalizacji.

\begin{lstlisting}
#include <concepts>
#include <string_view>

template<class T>
struct TypeLabel {
    static constexpr std::string_view value = "other";
};

template<std::integral T>
struct TypeLabel<T> {
    static constexpr std::string_view value = "integer";
};

template<std::floating_point T>
struct TypeLabel<T> {
    static constexpr std::string_view value = "floating";
};

static_assert(TypeLabel<int>::value == "integer");
static_assert(TypeLabel<double>::value == "floating");
\end{lstlisting}

Nie używaj jednak złożonego systemu specjalizacji, jeżeli
prosta funkcja przeciążona lub \texttt{if constexpr}
lepiej wyraża zamysł.

\section{Koncepcje i pakiety parametrów: kontynuacja rozdziału 4}
\label{sec:concepts-variadic}

Koncepcje bardzo dobrze współpracują z \emph{variadic templates}.
Możemy nakładać warunek na każdy element pakietu typów
za pomocą \emph{fold expression}.

\begin{lstlisting}
#include <concepts>

template<class... Ts>
requires (std::integral<Ts> && ...)
constexpr auto sum_integers(Ts... xs) {
    return (0 + ... + xs);
}

static_assert(sum_integers() == 0);
static_assert(sum_integers(1, 2, 3) == 6);
// sum_integers(1, 2.5);    // ERROR: double not integral
\end{lstlisting}

Zapis \texttt{(P<Ts> \&\& ...)} reprezentuje
koniunkcję wymagań dla wszystkich elementów pakietu.
Dla pustego pakietu taki fold logiczny daje \texttt{true};
nie oznacza to jednak, że każde ciało funkcji poradzi sobie
z zerową liczbą argumentów. W przykładzie użyliśmy
\texttt{0} jako elementu neutralnego sumowania.

\subsection{Ograniczony pakiet typów i uniwersalne przekazywanie}

\begin{lstlisting}
#include <concepts>
#include <utility>

template<std::integral... Ts>
void accept_integers(Ts&&... xs) {
    ((void)std::forward<Ts>(xs), ...); // illustrative use
}
\end{lstlisting}

Zapis \texttt{std::integral... Ts} również ogranicza wszystkie
elementy pakietu. Uwaga: przy dedukcji z przekazywanymi referencjami
\texttt{Ts\&\&}, typ \texttt{Ts} może być \emph{referencją},
gdy argument jest lvalue. W takim przypadku
\texttt{std::integral<Ts>} może być fałszywe.
Jeżeli chcemy akceptować obie kategorie wartości, często
właściwsze jest użycie \texttt{std::remove\_cvref\_t}:

\begin{lstlisting}
#include <concepts>
#include <type_traits>
#include <utility>

template<class... Ts>
requires (std::integral<std::remove_cvref_t<Ts>> && ...)
void accept_any_int(Ts&&... xs) {
    ((void)xs, ...); // demonstration; no ownership transfer
}
\end{lstlisting}

Ta różnica jest szczególnie ważna przy projektowaniu generycznych
wrapperów, które mają przyjmować także obiekty \texttt{const}
i referencje.

\section{Concepts a \texttt{if constexpr}, \texttt{static\_assert} i SFINAE}
\label{sec:concepts-porownanie}

Są to narzędzia pokrewne, ale mają różne cele.

\subsection{\texttt{if constexpr}: wybór implementacji}

\begin{lstlisting}
#include <concepts>
#include <string>
#include <string_view>

template<class T>
std::string describe(const T& value) {
    if constexpr (std::integral<T>) {
        return "an integer";
    } else if constexpr (std::floating_point<T>) {
        return "a floating-point value";
    } else {
        return "another type";
    }
}
\end{lstlisting}

\texttt{if constexpr} umożliwia wybranie gałęzi kodu
\emph{wewnątrz} szablonu. Nie tworzy samo przez się
ograniczonego interfejsu. Z kolei \texttt{requires}
może odrzucić niewłaściwe wywołanie jeszcze na etapie
dobierania przeciążeń.

\subsection{\texttt{static\_assert}: własny komunikat}

\begin{lstlisting}
#include <concepts>

template<class T>
void save_number(const T& value) {
    static_assert(std::integral<T>,
                  "save_number expects an integral type");
    // ...
}
\end{lstlisting}

\texttt{static\_assert} jest użyteczne, gdy potrzebujemy
jednoznacznego komunikatu albo sprawdzamy niezmiennik wewnątrz
implementacji. Nie zastępuje jednak ograniczenia uczestniczącego
w rozstrzyganiu przeciążeń.

\subsection{Przed C++20: \texttt{std::enable\_if} i idiom wykrywania}

\begin{lstlisting}
#include <type_traits>

// C++11/14/17 technique:
template<class T,
         std::enable_if_t<std::is_integral_v<T>, int> = 0>
T old_square(T x) {
    return x * x;
}

// C++20:
#include <concepts>
template<std::integral T>
T new_square(T x) {
    return x * x;
}
\end{lstlisting}

Dla nowego kodu C++20 najczęściej korzystniej wybrać
koncepcję: intencja staje się czytelna i spójna z mechanizmem
porządkowania przeciążeń. \texttt{std::enable\_if}
nadal ma znaczenie w bibliotekach wspierających starsze
wersje standardu.

\section{Praktyczny przykład: średnia z zakresu liczb}
\label{sec:concepts-praktyka}

Zamiast pisać funkcję tylko dla \texttt{std::vector<double>},
możemy określić dokładne wymagania wobec \emph{zakresu}.

\begin{lstlisting}
#include <concepts>
#include <cstddef>
#include <optional>
#include <ranges>
#include <vector>
#include <array>

template<class R>
concept NumericInputRange =
    std::ranges::input_range<R> &&
    std::convertible_to<std::ranges::range_reference_t<R>, double>;

template<NumericInputRange R>
std::optional<double> average(R&& range) {
    double sum = 0.0;
    std::size_t count = 0;
    for (auto&& item : range) {
        sum += static_cast<double>(item);
        ++count;
    }
    if (count == 0) return std::nullopt;
    return sum / static_cast<double>(count);
}

int main() {
    std::vector<int> v{2, 4, 6};
    std::array<double, 2> a{3.0, 7.0};
    auto x = average(v); // 4.0
    auto y = average(a); // 5.0
    auto z = average(std::vector<int>{}); // empty optional
    return !(x && y && !z && *x == 4.0 && *y == 5.0);
}
\end{lstlisting}

Funkcja działa dla różnych kontenerów i zakresów dostarczających
iterację jednokierunkową. Warunek \texttt{input\_range}
wyraża strukturę danych, a \texttt{convertible\_to}
-- oczekiwany rodzaj elementów. Zastosowaliśmy
\texttt{std::optional<double>}, aby pusty zakres nie powodował
dzielenia przez zero. Dla zastosowań numerycznych należy
dodatkowo rozważyć stabilność sumowania, przepełnienia i
znaczenie konwersji do \texttt{double}.

\paragraph{Kiedy \texttt{ranges}?}
Koncepcje standardowych zakresów są szczególnie użyteczne,
gdy funkcja operuje na sekwencji, ale nie wymaga konkretnie
\texttt{std::vector} czy \texttt{std::list}. Zamiast wymuszać
silniejsze wymaganie \texttt{random\_access\_range},
wybieraj najsłabszą koncepcję, która wystarcza do działania
algorytmu.

\section{Koncepcje dla lambd i algorytmów}

Funkcje anonimowe również mogą mieć ograniczone parametry.

\begin{lstlisting}
#include <concepts>

int main() {
    auto twice = [](std::integral auto x) { return x + x; };
    return twice(5) != 10;
}
\end{lstlisting}

Można też tworzyć generyczne algorytmy przyjmujące funkcje
spełniające \texttt{std::predicate}.

\begin{lstlisting}
#include <concepts>
#include <cstddef>
#include <vector>

template<class T, class Pred>
requires std::predicate<Pred&, const T&>
std::size_t count_if_custom(const std::vector<T>& values,
                            Pred pred) {
    std::size_t count = 0;
    for (const T& x : values)
        if (pred(x)) ++count;
    return count;
}
\end{lstlisting}

Zwróć uwagę na argumenty typu w
\texttt{std::predicate<Pred\&, const T\&>}:
sprawdzają możliwość wywołania \emph{konkretnego rodzaju}
predykatu z \emph{konkretnym rodzajem} argumentu.
Kategoria wartości ma znaczenie i nie należy bezrefleksyjnie
wstawiać tu samych nazw typów.

\section{Granice gwarancji: najczęstsze pułapki}
\label{sec:concepts-pulapki}

\subsection{Concept nie sprawdza dowolnej własności semantycznej}

\begin{lstlisting}
#include <concepts>

struct Strange {
    int value{};
    bool operator==(const Strange&) const { return false; }
};

static_assert(std::equality_comparable<Strange>);
// Syntactically valid, but equality lacks reflexivity.
\end{lstlisting}

Kompilator akceptuje składnię porównania, lecz
\texttt{a == a} jest zawsze fałszywe. Kod nie spełnia
matematycznych oczekiwań związanych z relacją równości.
To obowiązek autora typu, niezależny od samego mechanizmu
\texttt{concepts}.

\subsection{Koncepcje nie zamieniają błędów w poprawny kod}

Wyrażenia \texttt{requires} pozwalają w kontekście zależnym
od parametru szablonu traktować niepoprawne podstawienia
jako niespełnione wymagania. Nie oznacza to jednak, że
\emph{każdy} błędny fragment C++ staje się automatycznie
\texttt{false}. Niezależne od parametru, trwale niepoprawne
konstrukcje mogą prowadzić do błędu kompilacji.

\subsection{Różnica między \texttt{T}, \texttt{T\&} i \texttt{const T\&}}

\begin{lstlisting}
#include <concepts>
#include <type_traits>

template<class T>
concept PlainInteger =
    std::integral<std::remove_cvref_t<T>>;

static_assert(!std::integral<int&>);
static_assert(PlainInteger<int&>);
static_assert(PlainInteger<const int&>);
\end{lstlisting}

Przy dedukcji typu generycznego parametru
\texttt{T\&\&} ta różnica często decyduje o tym,
czy z pozoru poprawne wywołanie spełni koncepcję.

\subsection{\texttt{std::same\_as} czy \texttt{std::convertible\_to}?}

\begin{lstlisting}
#include <concepts>

template<class T>
concept StrictAdd = requires(T a, T b) {
    { a + b } -> std::same_as<T>;
};

template<class T>
concept FlexibleAdd = requires(T a, T b) {
    { a + b } -> std::convertible_to<T>;
};

static_assert(!StrictAdd<short>);  // short + short -> int
static_assert(FlexibleAdd<short>); // int -> short allowed
\end{lstlisting}

Koncepcja \texttt{FlexibleAdd} nie oznacza \emph{bezstratnej}
konwersji; \texttt{int} można niejawnie przekształcić do
\texttt{short}, tracąc informację. Nazwa koncepcji powinna
odzwierciedlać rzeczywistą, a nie postulowaną gwarancję.

\subsection{Funkcja ograniczona nie zawsze jest optymalniejsza}

Koncepcje opisują poprawność i wybór instancji szablonów.
Nie wymuszają \emph{inline}, wektoryzacji, eliminacji kopii
ani lepszego kodu maszynowego. Można je stosować bez obietnicy
przyspieszenia wykonywania programu.

\section{Zasady projektowania: kiedy używać, a kiedy nie?}

\begin{center}
\small
\begin{tabular}{|p{0.34\textwidth}|p{0.55\textwidth}|}
\hline
\textbf{Sytuacja} & \textbf{Zalecane rozwiązanie} \\
\hline
Funkcja generyczna wymaga konkretnej kategorii typów & Nazwany \texttt{concept}, np. \texttt{std::integral} \\
\hline
Wymagania dotyczą istnienia operacji lub typu składowego & Własny \texttt{requires-expression} \\
\hline
Ważne są różne wersje algorytmu dla różnych typów & Ograniczone przeciążenia lub specjalizacje \\
\hline
Jedna implementacja ma kilka gałęzi zależnych od typu & \texttt{if constexpr} \\
\hline
Sprawdzasz założenie wewnątrz ciała funkcji & \texttt{static\_assert} z komunikatem \\
\hline
Weryfikujesz jednakowy warunek dla całego pakietu & Fold \texttt{(Concept<Ts> \&\& ...)} \\
\hline
Potrzebujesz iteracji po sekwencji & \texttt{std::ranges::input\_range} lub silniejszy koncept \\
\hline
Funkcja działa dla jednego, konkretnego typu & Zwykła funkcja często jest prostsza niż szablon \\
\hline
Projekt musi działać także w C++17 & \texttt{type\_traits} i (w razie potrzeby) SFINAE \\
\hline
Wymagasz warunku wykonywania programu, np. \texttt{n>0} & Zwykłe sprawdzenie w czasie działania; koncept go nie zastępuje \\
\hline
\end{tabular}
\end{center}

\section{Podsumowanie i pytania kontrolne}

\begin{enumerate}
\item \texttt{concept} nazywa predykat nad parametrami szablonu,
sprawdzany w czasie kompilacji.
\item \texttt{requires-clause} ogranicza deklarację;
\texttt{requires-expression} jest wyrażeniem badającym wymagania.
\item Istnieją cztery formy wymagań: proste, typu, złożone
i zagnieżdżone.
\item \texttt{std::same\_as}, \texttt{std::convertible\_to},
\texttt{std::integral}, \texttt{std::predicate} i koncepcje zakresów
powinny być zwykle pierwszym wyborem.
\item Ograniczenia mogą wpływać na wybór przeciążenia, ale mechanizm
porządkowania nie jest dowolnym dowodzeniem implikacji logicznej.
\item Nazwane koncepcje warto wykorzystywać ponownie przy
budowaniu koncepcji bardziej szczegółowych.
\item \texttt{std::remove\_cvref\_t<T>} ma często znaczenie
przy koncepcjach typów dedukowanych z referencji przekazujących.
\item Koncepcje opisują przede wszystkim wymagania składniowe,
a dodatkowe prawa semantyczne muszą respektować autorzy typów.
\end{enumerate}

\paragraph{Pytania do samodzielnego sprawdzenia.}
Zastanów się, co kompilator zaakceptuje i dlaczego:

\begin{lstlisting}
#include <concepts>
#include <type_traits>

static_assert(std::integral<bool>);
static_assert(!std::integral<int&>);
static_assert(std::integral<std::remove_cvref_t<const int&>>);
static_assert(!std::same_as<int, long>);

// 1. Why does StrictAdd<short> fail?
// 2. Why can a concept be true for a logically broken ==?
// 3. How do requires-clause and requires-expression differ?
// 4. Which overload of category(5) is chosen?
// 5. Why is NumericInputRange more general than vector<double>?
\end{lstlisting}

\paragraph{Źródła i dalsza lektura.}
Oficjalne reguły języka określa standard ISO C++, a przystępne
opisy i przykłady znajdują się w serwisie \emph{cppreference}.
Przydatne adresy:
\begin{lstlisting}
en.cppreference.com/cpp/language/constraints
en.cppreference.com/cpp/language/requires
en.cppreference.com/cpp/concepts
\end{lstlisting}


\chapter[Concepts i requires w C++]{Koncepcje i ograniczenia szablonów (\textit{concepts} i \texttt{requires})}
\label{chap:concepts-cpp}

\section{Po co wprowadzono koncepcje?}

Szablony C++ pozwalają pisać kod niezależny od konkretnego typu.
Tradycyjny szablon deklaruje jednak zwykle \emph{co kompilator ma zrobić},
nie wyjaśniając wprost, \emph{jakich własności} wymaga od argumentów.
Jeśli operacja jest niepoprawna, błąd może zostać wykryty dopiero w ciele
szablonu, a komunikat diagnostyczny bywa długi.

\begin{lstlisting}
template<class T>
T twice(const T& x) {
    return x + x;           // requires operator+
}

// twice(7);                // OK
// twice(SomeType{});       // may fail inside function body
\end{lstlisting}

Od C++20 można umieścić wymagania na poziomie interfejsu:

\begin{lstlisting}
#include <concepts>

template<class T>
concept Addable = requires(const T& a, const T& b) {
    { a + b } -> std::convertible_to<T>;
};

template<Addable T>
T twice(const T& x) {
    return x + x;
}
\end{lstlisting}

Koncepcja (\textit{concept}) to nazwany predykat sprawdzany podczas
kompilacji, opisujący wymagania wobec argumentów szablonu.
Może dotyczyć istnienia operatorów, typów składowych, własności typów,
kategorii iteratorów i wielu innych cech. \emph{Constraints}
(ograniczenia) wpływają zarówno na dopuszczalność konkretyzacji,
jak i na wybór przeciążenia funkcji.

\paragraph{Wersja języka.}
W całym rozdziale zakładamy \textbf{C++20} oraz kompilację
np. poleceniem \texttt{g++ -std=c++20 -Wall -Wextra -Wpedantic main.cpp}.
Definicje koncepcji standardowych znajdują się przede wszystkim
w nagłówkach \texttt{<concepts>}, \texttt{<iterator>} i
\texttt{<ranges>}.

\section{Pierwsza koncepcja: składnia i zastosowanie}
\label{sec:concepts-pierwsza}

Definicja nazwanej koncepcji ma postać
\texttt{template<...> concept Nazwa = warunek;}. Warunek powinien
być wyrażeniem stałym typu \texttt{bool}. Definicja koncepcji
występuje na poziomie przestrzeni nazw; nie jest zwykłą funkcją.

\begin{lstlisting}
#include <concepts>
#include <type_traits>

template<class T>
concept Integer = std::is_integral_v<T>;

static_assert(Integer<int>);
static_assert(Integer<bool>);    // bool is integral
static_assert(!Integer<double>);
\end{lstlisting}

Koncepcję można sprawdzić wyrażeniem \texttt{Integer<int>} lub
\texttt{static\_assert}, ale nie należy jej mylić ze zwykłą
zmienną. W szczególności \texttt{concept} nie daje się jawnie
specjalizować w sposób analogiczny do szablonu klasy.

\subsection{Cztery podstawowe sposoby ograniczania funkcji}

Poniższe deklaracje opisują ten sam warunek: funkcja jest przeznaczona
dla typów całkowitych w sensie \texttt{std::integral}.

\begin{lstlisting}
#include <concepts>

template<class T>
requires std::integral<T>
T square1(T x) { return x * x; }

template<class T>
T square2(T x) requires std::integral<T> {
    return x * x;
}

template<std::integral T>
T square3(T x) { return x * x; }

std::integral auto square4(std::integral auto x) {
    return x * x;
}
\end{lstlisting}

W ostatnim przykładzie ograniczony jest nie tylko typ argumentu,
ale też \emph{typ wyniku}. Sama dedukcja typu wyniku nadal odbywa
się na podstawie wyrażenia \texttt{x * x}; koncept jest sprawdzany
po jej przeprowadzeniu. Zapis \texttt{std::integral auto} w parametrze
jest skrótem szablonu funkcji, tzw. \emph{abbreviated function template}.

\paragraph{Kiedy którą postać wybrać?}
Dla jednego prostego ograniczenia zapis z ograniczonym
parametrem szablonu, np. \texttt{template<\allowbreak std::integral T>},
jest krótki i czytelny. Jeśli zależność dotyczy kilku parametrów,
najwygodniejsza bywa klauzula \texttt{requires}.
Skrócona postać z \texttt{auto} przydaje się w niewielkich funkcjach
i lambdach.

\section{Dwie role słowa \texttt{requires}}

\begin{center}\small
\begin{tabularx}{\textwidth}{@{}p{0.34\textwidth}X@{}}
\toprule
Zapis & Za co odpowiada \\
\midrule
\code{concept C = ...} & Nadaje nazwę logicznemu ograniczeniu parametrów szablonu. \\
\code{requires C<T>} & Ogranicza udział szablonu w przeciążaniu do typów spełniających \code{C}. \\
\code{requires(T x) { x.f(); }} & Sprawdza poprawność wskazanego wyrażenia dla podstawionego typu. \\
\code{requires requires(...) { ... }} & Pierwsze \code{requires} wprowadza ograniczenie, drugie wyrażenie sprawdzające. \\
\code{static_assert(C<T>)} & Wymusza spełnienie warunku w danym miejscu i zgłasza błąd kompilacji, gdy jest fałszywy. \\
\bottomrule
\end{tabularx}
\end{center}
\label{sec:concepts-dwa-requires}

Identyczny zapis słowny \texttt{requires} występuje w dwóch
różnych konstrukcjach:
\begin{enumerate}
  \item \emph{requires-clause} -- klauzula ograniczająca deklarację;
  \item \emph{requires-expression} -- wyrażenie logiczne, które bada
  poprawność podanych wymagań.
\end{enumerate}

\begin{lstlisting}
#include <concepts>

template<class T>
concept Incrementable = requires(T value) { // requires-expression
    ++value;
};

template<class T>
requires Incrementable<T>                // requires-clause
void increment(T& x) {
    ++x;
}
\end{lstlisting}

Można je połączyć bez nadawania koncepcji osobnej nazwy:

\begin{lstlisting}
template<class T>
requires requires(T x) { ++x; }
void increment_again(T& x) {
    ++x;
}
\end{lstlisting}

Pierwsze \texttt{requires} w ostatnim przykładzie zaczyna klauzulę;
drugie wprowadza wyrażenie sprawdzające wymagania. To poprawna
składnia, choć przy wielokrotnym użyciu lepiej utworzyć nazwaną
koncepcję.

\section{Cztery rodzaje wymagań w \texttt{requires-expression}}

\begin{center}\small
\begin{tabularx}{\textwidth}{@{}p{0.31\textwidth}X@{}}
\toprule
Rodzaj wymagania & Co sprawdza \\
\midrule
Proste: \code{a + b;} & Czy wyrażenie jest poprawne składniowo dla danych typów. \\
Typowe: \code{typename T::value_type;} & Czy wskazany typ zależny istnieje. \\
Złożone: \code{{ f(x) } -> C;} & Czy wyrażenie jest poprawne oraz czy jego typ spełnia ograniczenie \code{C}. \\
Zagnieżdżone: \code{requires C<T>;} & Czy dodatkowy warunek logiczny jest spełniony. \\
\bottomrule
\end{tabularx}
\end{center}
\label{sec:concepts-cztery-wymagania}

Wyrażenie \texttt{requires} umożliwia opisanie czterech rodzajów
warunków. Parametry wpisane w nawiasach \texttt{requires(T x)}
są jedynie \emph{lokalnymi parametrami do sprawdzania wyrażeń};
nie tworzy się dla nich obiektów podczas działania programu.

\subsection{Wymaganie proste (\textit{simple requirement})}

\begin{lstlisting}
#include <ostream>

template<class T>
concept Printable = requires(std::ostream& os, const T& x) {
    os << x;                    // must be well-formed
};
\end{lstlisting}

Wyrażenie \texttt{os << x}
nie jest tutaj wykonywane; kompilator sprawdza jego poprawność.
Proste wymaganie nie nakłada warunku na typ zwracany.

\subsection{Wymaganie typu (\textit{type requirement})}

\begin{lstlisting}
#include <vector>

template<class T>
concept HasValueType = requires {
    typename T::value_type;    // nested type must exist
};

static_assert(HasValueType<std::vector<int>>);
static_assert(!HasValueType<int>);
\end{lstlisting}

Słowo \texttt{typename} jest tu składnikiem składni wymagania,
a nie deklaracją nowego typu.

\subsection{Wymaganie złożone (\textit{compound requirement})}

\begin{lstlisting}
#include <concepts>

template<class T>
concept AddableNoThrow = requires(const T& a, const T& b) {
    { a + b } noexcept -> std::same_as<T>;
};
\end{lstlisting}

Zapis \texttt{\{ wyrazenie \} noexcept -> Koncepcja;}
może sprawdzać jednocześnie: poprawność wyrażenia,
to, że jest ono \texttt{noexcept}, i własność typu
\texttt{decltype((wyrazenie))}. Obie części, \texttt{noexcept}
oraz \texttt{-> Koncepcja}, są opcjonalne.

\begin{lstlisting}
template<class T>
concept PlusReturnsT = requires(const T& a, const T& b) {
    { a + b } -> std::same_as<T>;
};

static_assert(PlusReturnsT<int>);
// A type with a proxy return value may not satisfy this.
\end{lstlisting}

\paragraph{Uwaga na typ i konwersje.}
Koncepcja \texttt{std::same\_as<T>} żąda \emph{dokładnie}
typu \texttt{T}, natomiast koncepcja
\texttt{std::convertible\_\allowbreak to<T>} dopuszcza konwersję.
Dla \texttt{int} operator \texttt{+} daje \texttt{int}, ale np.
dla małych typów całkowitych może zajść promocja do \texttt{int}.

\subsection{Wymaganie zagnieżdżone (\textit{nested requirement})}

\begin{lstlisting}
#include <type_traits>

template<class T>
concept SmallTrivial = requires {
    requires sizeof(T) <= 16;
    requires std::is_trivially_copyable_v<T>;
};

static_assert(SmallTrivial<int>);
\end{lstlisting}

Wymaganie zagnieżdżone zaczyna się kolejnym \texttt{requires}
i musi zawierać warunek stały. Jest wygodne, gdy wymaganie
dotyczy własności typu, a nie określonego wyrażenia.

\subsection{Przykład: połączenie różnych wymagań}

\begin{lstlisting}
#include <concepts>
#include <cstddef>

template<class T>
concept Indexable = requires(T& obj, std::size_t i) {
    typename T::value_type;
    { obj.size() } -> std::convertible_to<std::size_t>;
    obj[i];
};
\end{lstlisting}

To przykład \emph{minimalnego kontraktu składniowego}.
Koncepcja nie dowodzi, że indeks \texttt{i} jest w zakresie,
że operacja jest szybka ani że wynik indeksowania zachowuje się
jak zwykła referencja. Takie własności wymagają dodatkowych założeń.

\section{Standardowe koncepcje z biblioteki C++20}
\label{sec:concepts-standardowe}

Przed napisaniem własnej koncepcji warto sprawdzić, czy istnieje
gotowa w standardowej bibliotece.

\subsection{Typy, konwersje i dziedziczenie}

\begin{lstlisting}
#include <concepts>

static_assert(std::same_as<int, int>);
static_assert(!std::same_as<int, long>);
static_assert(std::integral<unsigned>);
static_assert(std::signed_integral<int>);
static_assert(std::unsigned_integral<unsigned>);
static_assert(std::floating_point<double>);
static_assert(std::convertible_to<int, double>);

struct Base {};
struct Derived : Base {};
static_assert(std::derived_from<Derived, Base>);
\end{lstlisting}

\texttt{std::derived\_from<D,B>} obejmuje publiczne i jednoznaczne
dziedziczenie oraz przypadek tego samego typu. Różni się pod tym
względem od samego \texttt{std::is\_base\_of\_v}.

\subsection{Konstruowalność, kopiowanie, przenoszenie i porównywanie}

Przydatne są m.in. koncepcje wyrażające inicjalizowalność,
konstruowalność, możliwość przenoszenia i kopiowania oraz porównywania:
\begin{lstlisting}
default_initializable    constructible_from
move_constructible       copy_constructible
movable                  copyable
semiregular              regular
equality_comparable      totally_ordered
\end{lstlisting}

\begin{lstlisting}
#include <concepts>
#include <memory>
#include <string>

static_assert(std::regular<int>);
static_assert(std::copyable<std::string>);
static_assert(std::movable<std::unique_ptr<int>>);
static_assert(!std::copyable<std::unique_ptr<int>>);
static_assert(std::totally_ordered<int>);
\end{lstlisting}

Koncepcje te opisują interfejs i zakładają także określone
\emph{prawa semantyczne} (np. sensowne porównywanie). Sam kompilator
nie potrafi udowodnić dowolnej własności semantycznej implementacji.

\subsection{Funkcje wywoływalne i predykaty}

\begin{lstlisting}
#include <concepts>

struct IsPositive {
    bool operator()(int x) const { return x > 0; }
};

static_assert(std::invocable<IsPositive, int>);
static_assert(std::predicate<IsPositive, int>);
static_assert(!std::predicate<IsPositive, void*>);
\end{lstlisting}

\texttt{std::invocable<F,Args...>} sprawdza możliwość wywołania,
natomiast \texttt{std::predicate} wymaga dodatkowo wyniku,
którego można użyć jako warunku logicznego, i wiąże się
z określonymi wymaganiami semantycznymi.

\subsection{Iteratory i zakresy}

Koncepcje iteratorów są udostępniane przez \texttt{<iterator>},
a koncepcje zakresów przez \texttt{<ranges>}.

\begin{lstlisting}
#include <ranges>
#include <vector>
#include <list>

static_assert(std::ranges::random_access_range<
              std::vector<int>>);
static_assert(std::ranges::bidirectional_range<
              std::list<int>>);
static_assert(!std::ranges::random_access_range<
              std::list<int>>);
\end{lstlisting}

W szczególności przydatne są koncepcje \emph{input},
\emph{forward}, \emph{bidirectional}, \emph{random access}
i \emph{contiguous range}. Pozwalają opisywać wymagania
algorytmów bez uzależniania ich od konkretnego kontenera.

\section{Łączenie koncepcji i definiowanie własnych kontraktów}

Wymagania można łączyć operatorami logicznymi
\texttt{\&\&} i \texttt{||}. Często lepiej budować niewielkie,
nazwane koncepcje niż powtarzać jedno długie wyrażenie.

\begin{lstlisting}
#include <concepts>

template<class T>
concept Arithmetic = std::integral<T> || std::floating_point<T>;

template<class T>
concept NonBooleanInteger =
    std::integral<T> && !std::same_as<T, bool>;

template<Arithmetic T>
T cube(T x) { return x * x * x; }

static_assert(Arithmetic<double>);
static_assert(!NonBooleanInteger<bool>);
\end{lstlisting}

Dla połączenia \texttt{A \&\& B} wymaganie \texttt{B} jest
sprawdzane dopiero, gdy \texttt{A} jest spełnione; analogicznie
dla \texttt{A || B} drugi składnik jest potrzebny dopiero,
gdy pierwszy nie jest spełniony. Ta reguła chroni często przed
podstawianiem wyrażeń, które dla danego typu nie mają sensu.

\begin{lstlisting}
#include <concepts>
#include <type_traits>

template<class T>
concept HasIntegralValueType =
    requires { typename T::value_type; } &&
    std::integral<typename T::value_type>;
\end{lstlisting}

W powyższym przykładzie odwołanie do \texttt{T::value\_type}
po prawej stronie jest istotne dopiero po potwierdzeniu,
że typ składowy istnieje. Należy jednak pamiętać, że złożone
reguły podstawiania i normalizacji koncepcji mają swoje szczegóły;
nie każdą źle sformułowaną deklarację naprawi samo zwarcie logiczne.

\section{Koncepcje z wieloma parametrami}

Koncepcja nie musi opisywać własności jednego typu.
Może także opisywać \emph{relację między typami}.

\begin{lstlisting}
#include <concepts>

template<class From, class To>
concept SameType = std::same_as<From, To>;

template<class Derived, class Base>
concept PublicChildOf = std::derived_from<Derived, Base>;

template<class T, class U>
concept AddableTogether = requires(const T& a, const U& b) {
    a + b;
};

static_assert(AddableTogether<int, double>);
\end{lstlisting}

Wymienione koncepcje \texttt{SameType} i
\texttt{PublicChildOf} służą tu jedynie ilustracji składni.
W praktyce używaj bezpośrednio standardowych
\texttt{std::same\_as} i \texttt{std::derived\_from},
jeżeli nie dodajesz do nich istotnego znaczenia domenowego.

\paragraph{Wstawianie argumentu przy ograniczonym parametrze.}

\begin{lstlisting}
struct Animal {};

template<std::derived_from<Animal> T>
void inspect(const T&) {}

// Equivalent constraint: std::derived_from<T, Animal>
\end{lstlisting}

W zapisie powyżej typ \texttt{T} jest podstawiany
\emph{na pierwszej pozycji} koncepcji, zatem
chodzi o warunek z poniższego przykładu:
\begin{lstlisting}
std::derived_from<T, Animal>
\end{lstlisting}

\section{Przeciążanie funkcji i wybór bardziej szczegółowego ograniczenia}
\label{sec:concepts-przeciazanie}

Ograniczenia uczestniczą w wyborze przeciążenia. Jeśli dwie funkcje
są dopasowane do argumentów, kompilator może preferować tę,
której ograniczenie jest bardziej szczegółowe (ang.
\emph{subsumption}).

\begin{lstlisting}
#include <concepts>
#include <string_view>

template<class T>
requires std::integral<T>
constexpr std::string_view category(T) {
    return "integral";
}

template<class T>
requires std::signed_integral<T>
constexpr std::string_view category(T) {
    return "signed integral";
}

static_assert(category(12) == "signed integral");
static_assert(category(12u) == "integral");
\end{lstlisting}

\texttt{std::signed\_integral} jest zbudowane na
\texttt{std::integral} i dodaje dodatkowe wymaganie. Dzięki temu
bardziej szczegółowe przeciążenie zostaje wybrane dla \texttt{int}.

\subsection{Ważna pułapka: podobny zapis to nie to samo ograniczenie}

Porządek ograniczeń nie jest ogólnym dowodzeniem implikacji
matematycznych. Kompilator porównuje znormalizowane ograniczenia
atomowe i ich \emph{tożsamość}, a nie tylko sens pozornie podobnych
wyrażeń.

\begin{lstlisting}
#include <type_traits>

template<class T>
concept HasSize = requires(const T& x) { x.size(); };

template<class T>
concept HasSizeAndCopy =
    HasSize<T> && std::is_copy_constructible_v<T>;

template<class T> requires HasSize<T>
int classify(const T&) { return 1; }

template<class T> requires HasSizeAndCopy<T>
int classify(const T&) { return 2; }
\end{lstlisting}

Druga koncepcja \emph{wykorzystuje tę samą} koncepcję
\texttt{HasSize}, dlatego kompilator rozpoznaje wspólny
składnik. Jeśli powtórzymy ręcznie nawet identycznie wyglądające
wymaganie opisujące metodę \texttt{size()},
nie mamy automatycznej gwarancji identyczności ograniczeń
atomowych.

\paragraph{Rekomendacja.}
Twórz proste nazwane koncepcje i buduj bardziej szczegółowe
z ich wykorzystaniem. Ułatwia to konserwację kodu i
przewidywanie wyboru przeciążeń.

\section{Koncepcje dla szablonów klas i ich specjalizacji}

Ograniczać można nie tylko funkcje, lecz również szablony klas.

\begin{lstlisting}
#include <concepts>

template<std::integral T>
class Counter {
    T value_{};
public:
    void add(T delta) { value_ += delta; }
    T value() const { return value_; }
};

// Counter<int> a;       // OK
// Counter<double> b;    // ERROR: constraint not satisfied
\end{lstlisting}

Ograniczenia mogą też sterować wyborem częściowej specjalizacji.

\begin{lstlisting}
#include <concepts>
#include <string_view>

template<class T>
struct TypeLabel {
    static constexpr std::string_view value = "other";
};

template<std::integral T>
struct TypeLabel<T> {
    static constexpr std::string_view value = "integer";
};

template<std::floating_point T>
struct TypeLabel<T> {
    static constexpr std::string_view value = "floating";
};

static_assert(TypeLabel<int>::value == "integer");
static_assert(TypeLabel<double>::value == "floating");
\end{lstlisting}

Nie używaj jednak złożonego systemu specjalizacji, jeżeli
prosta funkcja przeciążona lub \texttt{if constexpr}
lepiej wyraża zamysł.

\section{Koncepcje i pakiety parametrów: kontynuacja rozdziału 4}
\label{sec:concepts-variadic}

Koncepcje bardzo dobrze współpracują z \emph{variadic templates}.
Możemy nakładać warunek na każdy element pakietu typów
za pomocą \emph{fold expression}.

\begin{lstlisting}
#include <concepts>

template<class... Ts>
requires (std::integral<Ts> && ...)
constexpr auto sum_integers(Ts... xs) {
    return (0 + ... + xs);
}

static_assert(sum_integers() == 0);
static_assert(sum_integers(1, 2, 3) == 6);
// sum_integers(1, 2.5);    // ERROR: double not integral
\end{lstlisting}

Zapis \texttt{(P<Ts> \&\& ...)} reprezentuje
koniunkcję wymagań dla wszystkich elementów pakietu.
Dla pustego pakietu taki fold logiczny daje \texttt{true};
nie oznacza to jednak, że każde ciało funkcji poradzi sobie
z zerową liczbą argumentów. W przykładzie użyliśmy
\texttt{0} jako elementu neutralnego sumowania.

\subsection{Ograniczony pakiet typów i uniwersalne przekazywanie}

\begin{lstlisting}
#include <concepts>
#include <utility>

template<std::integral... Ts>
void accept_integers(Ts&&... xs) {
    ((void)std::forward<Ts>(xs), ...); // illustrative use
}
\end{lstlisting}

Zapis \texttt{std::integral... Ts} również ogranicza wszystkie
elementy pakietu. Uwaga: przy dedukcji z przekazywanymi referencjami
\texttt{Ts\&\&}, typ \texttt{Ts} może być \emph{referencją},
gdy argument jest lvalue. W takim przypadku
\texttt{std::integral<Ts>} może być fałszywe.
Jeżeli chcemy akceptować obie kategorie wartości, często
właściwsze jest użycie \texttt{std::remove\_cvref\_t}:

\begin{lstlisting}
#include <concepts>
#include <type_traits>
#include <utility>

template<class... Ts>
requires (std::integral<std::remove_cvref_t<Ts>> && ...)
void accept_any_int(Ts&&... xs) {
    ((void)xs, ...); // demonstration; no ownership transfer
}
\end{lstlisting}

Ta różnica jest szczególnie ważna przy projektowaniu generycznych
wrapperów, które mają przyjmować także obiekty \texttt{const}
i referencje.

\section{Concepts a \texttt{if constexpr}, \texttt{static\_assert} i SFINAE}
\label{sec:concepts-porownanie}

Są to narzędzia pokrewne, ale mają różne cele.

\subsection{\texttt{if constexpr}: wybór implementacji}

\begin{lstlisting}
#include <concepts>
#include <string>
#include <string_view>

template<class T>
std::string describe(const T& value) {
    if constexpr (std::integral<T>) {
        return "an integer";
    } else if constexpr (std::floating_point<T>) {
        return "a floating-point value";
    } else {
        return "another type";
    }
}
\end{lstlisting}

\texttt{if constexpr} umożliwia wybranie gałęzi kodu
\emph{wewnątrz} szablonu. Nie tworzy samo przez się
ograniczonego interfejsu. Z kolei \texttt{requires}
może odrzucić niewłaściwe wywołanie jeszcze na etapie
dobierania przeciążeń.

\subsection{\texttt{static\_assert}: własny komunikat}

\begin{lstlisting}
#include <concepts>

template<class T>
void save_number(const T& value) {
    static_assert(std::integral<T>,
                  "save_number expects an integral type");
    // ...
}
\end{lstlisting}

\texttt{static\_assert} jest użyteczne, gdy potrzebujemy
jednoznacznego komunikatu albo sprawdzamy niezmiennik wewnątrz
implementacji. Nie zastępuje jednak ograniczenia uczestniczącego
w rozstrzyganiu przeciążeń.

\subsection{Przed C++20: \texttt{std::enable\_if} i idiom wykrywania}

\begin{lstlisting}
#include <type_traits>

// C++11/14/17 technique:
template<class T,
         std::enable_if_t<std::is_integral_v<T>, int> = 0>
T old_square(T x) {
    return x * x;
}

// C++20:
#include <concepts>
template<std::integral T>
T new_square(T x) {
    return x * x;
}
\end{lstlisting}

Dla nowego kodu C++20 najczęściej korzystniej wybrać
koncepcję: intencja staje się czytelna i spójna z mechanizmem
porządkowania przeciążeń. \texttt{std::enable\_if}
nadal ma znaczenie w bibliotekach wspierających starsze
wersje standardu.

\section{Praktyczny przykład: średnia z zakresu liczb}
\label{sec:concepts-praktyka}

Zamiast pisać funkcję tylko dla \texttt{std::vector<double>},
możemy określić dokładne wymagania wobec \emph{zakresu}.

\begin{lstlisting}
#include <concepts>
#include <cstddef>
#include <optional>
#include <ranges>
#include <vector>
#include <array>

template<class R>
concept NumericInputRange =
    std::ranges::input_range<R> &&
    std::convertible_to<std::ranges::range_reference_t<R>, double>;

template<NumericInputRange R>
std::optional<double> average(R&& range) {
    double sum = 0.0;
    std::size_t count = 0;
    for (auto&& item : range) {
        sum += static_cast<double>(item);
        ++count;
    }
    if (count == 0) return std::nullopt;
    return sum / static_cast<double>(count);
}

int main() {
    std::vector<int> v{2, 4, 6};
    std::array<double, 2> a{3.0, 7.0};
    auto x = average(v); // 4.0
    auto y = average(a); // 5.0
    auto z = average(std::vector<int>{}); // empty optional
    return !(x && y && !z && *x == 4.0 && *y == 5.0);
}
\end{lstlisting}

Funkcja działa dla różnych kontenerów i zakresów dostarczających
iterację jednokierunkową. Warunek \texttt{input\_range}
wyraża strukturę danych, a \texttt{convertible\_to}
-- oczekiwany rodzaj elementów. Zastosowaliśmy
\texttt{std::optional<double>}, aby pusty zakres nie powodował
dzielenia przez zero. Dla zastosowań numerycznych należy
dodatkowo rozważyć stabilność sumowania, przepełnienia i
znaczenie konwersji do \texttt{double}.

\paragraph{Kiedy \texttt{ranges}?}
Koncepcje standardowych zakresów są szczególnie użyteczne,
gdy funkcja operuje na sekwencji, ale nie wymaga konkretnie
\texttt{std::vector} czy \texttt{std::list}. Zamiast wymuszać
silniejsze wymaganie \texttt{random\_access\_range},
wybieraj najsłabszą koncepcję, która wystarcza do działania
algorytmu.

\section{Koncepcje dla lambd i algorytmów}

Funkcje anonimowe również mogą mieć ograniczone parametry.

\begin{lstlisting}
#include <concepts>

int main() {
    auto twice = [](std::integral auto x) { return x + x; };
    return twice(5) != 10;
}
\end{lstlisting}

Można też tworzyć generyczne algorytmy przyjmujące funkcje
spełniające \texttt{std::predicate}.

\begin{lstlisting}
#include <concepts>
#include <cstddef>
#include <vector>

template<class T, class Pred>
requires std::predicate<Pred&, const T&>
std::size_t count_if_custom(const std::vector<T>& values,
                            Pred pred) {
    std::size_t count = 0;
    for (const T& x : values)
        if (pred(x)) ++count;
    return count;
}
\end{lstlisting}

Zwróć uwagę na argumenty typu w
\texttt{std::predicate<Pred\&, const T\&>}:
sprawdzają możliwość wywołania \emph{konkretnego rodzaju}
predykatu z \emph{konkretnym rodzajem} argumentu.
Kategoria wartości ma znaczenie i nie należy bezrefleksyjnie
wstawiać tu samych nazw typów.

\section{Granice gwarancji: najczęstsze pułapki}
\label{sec:concepts-pulapki}

\subsection{Concept nie sprawdza dowolnej własności semantycznej}

\begin{lstlisting}
#include <concepts>

struct Strange {
    int value{};
    bool operator==(const Strange&) const { return false; }
};

static_assert(std::equality_comparable<Strange>);
// Syntactically valid, but equality lacks reflexivity.
\end{lstlisting}

Kompilator akceptuje składnię porównania, lecz
\texttt{a == a} jest zawsze fałszywe. Kod nie spełnia
matematycznych oczekiwań związanych z relacją równości.
To obowiązek autora typu, niezależny od samego mechanizmu
\texttt{concepts}.

\subsection{Koncepcje nie zamieniają błędów w poprawny kod}

Wyrażenia \texttt{requires} pozwalają w kontekście zależnym
od parametru szablonu traktować niepoprawne podstawienia
jako niespełnione wymagania. Nie oznacza to jednak, że
\emph{każdy} błędny fragment C++ staje się automatycznie
\texttt{false}. Niezależne od parametru, trwale niepoprawne
konstrukcje mogą prowadzić do błędu kompilacji.

\subsection{Różnica między \texttt{T}, \texttt{T\&} i \texttt{const T\&}}

\begin{lstlisting}
#include <concepts>
#include <type_traits>

template<class T>
concept PlainInteger =
    std::integral<std::remove_cvref_t<T>>;

static_assert(!std::integral<int&>);
static_assert(PlainInteger<int&>);
static_assert(PlainInteger<const int&>);
\end{lstlisting}

Przy dedukcji typu generycznego parametru
\texttt{T\&\&} ta różnica często decyduje o tym,
czy z pozoru poprawne wywołanie spełni koncepcję.

\subsection{\texttt{std::same\_as} czy \texttt{std::convertible\_to}?}

\begin{lstlisting}
#include <concepts>

template<class T>
concept StrictAdd = requires(T a, T b) {
    { a + b } -> std::same_as<T>;
};

template<class T>
concept FlexibleAdd = requires(T a, T b) {
    { a + b } -> std::convertible_to<T>;
};

static_assert(!StrictAdd<short>);  // short + short -> int
static_assert(FlexibleAdd<short>); // int -> short allowed
\end{lstlisting}

Koncepcja \texttt{FlexibleAdd} nie oznacza \emph{bezstratnej}
konwersji; \texttt{int} można niejawnie przekształcić do
\texttt{short}, tracąc informację. Nazwa koncepcji powinna
odzwierciedlać rzeczywistą, a nie postulowaną gwarancję.

\subsection{Funkcja ograniczona nie zawsze jest optymalniejsza}

Koncepcje opisują poprawność i wybór instancji szablonów.
Nie wymuszają \emph{inline}, wektoryzacji, eliminacji kopii
ani lepszego kodu maszynowego. Można je stosować bez obietnicy
przyspieszenia wykonywania programu.

\section{Zasady projektowania: kiedy używać, a kiedy nie?}

\begin{center}
\small
\begin{tabular}{|p{0.34\textwidth}|p{0.55\textwidth}|}
\hline
\textbf{Sytuacja} & \textbf{Zalecane rozwiązanie} \\
\hline
Funkcja generyczna wymaga konkretnej kategorii typów & Nazwany \texttt{concept}, np. \texttt{std::integral} \\
\hline
Wymagania dotyczą istnienia operacji lub typu składowego & Własny \texttt{requires-expression} \\
\hline
Ważne są różne wersje algorytmu dla różnych typów & Ograniczone przeciążenia lub specjalizacje \\
\hline
Jedna implementacja ma kilka gałęzi zależnych od typu & \texttt{if constexpr} \\
\hline
Sprawdzasz założenie wewnątrz ciała funkcji & \texttt{static\_assert} z komunikatem \\
\hline
Weryfikujesz jednakowy warunek dla całego pakietu & Fold \texttt{(Concept<Ts> \&\& ...)} \\
\hline
Potrzebujesz iteracji po sekwencji & \texttt{std::ranges::input\_range} lub silniejszy koncept \\
\hline
Funkcja działa dla jednego, konkretnego typu & Zwykła funkcja często jest prostsza niż szablon \\
\hline
Projekt musi działać także w C++17 & \texttt{type\_traits} i (w razie potrzeby) SFINAE \\
\hline
Wymagasz warunku wykonywania programu, np. \texttt{n>0} & Zwykłe sprawdzenie w czasie działania; koncept go nie zastępuje \\
\hline
\end{tabular}
\end{center}

\section{Podsumowanie i pytania kontrolne}

\begin{enumerate}
\item \texttt{concept} nazywa predykat nad parametrami szablonu,
sprawdzany w czasie kompilacji.
\item \texttt{requires-clause} ogranicza deklarację;
\texttt{requires-expression} jest wyrażeniem badającym wymagania.
\item Istnieją cztery formy wymagań: proste, typu, złożone
i zagnieżdżone.
\item \texttt{std::same\_as}, \texttt{std::convertible\_to},
\texttt{std::integral}, \texttt{std::predicate} i koncepcje zakresów
powinny być zwykle pierwszym wyborem.
\item Ograniczenia mogą wpływać na wybór przeciążenia, ale mechanizm
porządkowania nie jest dowolnym dowodzeniem implikacji logicznej.
\item Nazwane koncepcje warto wykorzystywać ponownie przy
budowaniu koncepcji bardziej szczegółowych.
\item \texttt{std::remove\_cvref\_t<T>} ma często znaczenie
przy koncepcjach typów dedukowanych z referencji przekazujących.
\item Koncepcje opisują przede wszystkim wymagania składniowe,
a dodatkowe prawa semantyczne muszą respektować autorzy typów.
\end{enumerate}

\paragraph{Pytania do samodzielnego sprawdzenia.}
Zastanów się, co kompilator zaakceptuje i dlaczego:

\begin{lstlisting}
#include <concepts>
#include <type_traits>

static_assert(std::integral<bool>);
static_assert(!std::integral<int&>);
static_assert(std::integral<std::remove_cvref_t<const int&>>);
static_assert(!std::same_as<int, long>);

// 1. Why does StrictAdd<short> fail?
// 2. Why can a concept be true for a logically broken ==?
// 3. How do requires-clause and requires-expression differ?
// 4. Which overload of category(5) is chosen?
// 5. Why is NumericInputRange more general than vector<double>?
\end{lstlisting}

\paragraph{Źródła i dalsza lektura.}
Oficjalne reguły języka określa standard ISO C++, a przystępne
opisy i przykłady znajdują się w serwisie \emph{cppreference}.
Przydatne adresy:
\begin{lstlisting}
en.cppreference.com/cpp/language/constraints
en.cppreference.com/cpp/language/requires
en.cppreference.com/cpp/concepts
\end{lstlisting}


\chapter[Concepts i requires w C++]{Koncepcje i ograniczenia szablonów (\textit{concepts} i \texttt{requires})}
\label{chap:concepts-cpp}

\section{Po co wprowadzono koncepcje?}

Szablony C++ pozwalają pisać kod niezależny od konkretnego typu.
Tradycyjny szablon deklaruje jednak zwykle \emph{co kompilator ma zrobić},
nie wyjaśniając wprost, \emph{jakich własności} wymaga od argumentów.
Jeśli operacja jest niepoprawna, błąd może zostać wykryty dopiero w ciele
szablonu, a komunikat diagnostyczny bywa długi.

\begin{lstlisting}
template<class T>
T twice(const T& x) {
    return x + x;           // requires operator+
}

// twice(7);                // OK
// twice(SomeType{});       // may fail inside function body
\end{lstlisting}

Od C++20 można umieścić wymagania na poziomie interfejsu:

\begin{lstlisting}
#include <concepts>

template<class T>
concept Addable = requires(const T& a, const T& b) {
    { a + b } -> std::convertible_to<T>;
};

template<Addable T>
T twice(const T& x) {
    return x + x;
}
\end{lstlisting}

Koncepcja (\textit{concept}) to nazwany predykat sprawdzany podczas
kompilacji, opisujący wymagania wobec argumentów szablonu.
Może dotyczyć istnienia operatorów, typów składowych, własności typów,
kategorii iteratorów i wielu innych cech. \emph{Constraints}
(ograniczenia) wpływają zarówno na dopuszczalność konkretyzacji,
jak i na wybór przeciążenia funkcji.

\paragraph{Wersja języka.}
W całym rozdziale zakładamy \textbf{C++20} oraz kompilację
np. poleceniem \texttt{g++ -std=c++20 -Wall -Wextra -Wpedantic main.cpp}.
Definicje koncepcji standardowych znajdują się przede wszystkim
w nagłówkach \texttt{<concepts>}, \texttt{<iterator>} i
\texttt{<ranges>}.

\section{Pierwsza koncepcja: składnia i zastosowanie}
\label{sec:concepts-pierwsza}

Definicja nazwanej koncepcji ma postać
\texttt{template<...> concept Nazwa = warunek;}. Warunek powinien
być wyrażeniem stałym typu \texttt{bool}. Definicja koncepcji
występuje na poziomie przestrzeni nazw; nie jest zwykłą funkcją.

\begin{lstlisting}
#include <concepts>
#include <type_traits>

template<class T>
concept Integer = std::is_integral_v<T>;

static_assert(Integer<int>);
static_assert(Integer<bool>);    // bool is integral
static_assert(!Integer<double>);
\end{lstlisting}

Koncepcję można sprawdzić wyrażeniem \texttt{Integer<int>} lub
\texttt{static\_assert}, ale nie należy jej mylić ze zwykłą
zmienną. W szczególności \texttt{concept} nie daje się jawnie
specjalizować w sposób analogiczny do szablonu klasy.

\subsection{Cztery podstawowe sposoby ograniczania funkcji}

Poniższe deklaracje opisują ten sam warunek: funkcja jest przeznaczona
dla typów całkowitych w sensie \texttt{std::integral}.

\begin{lstlisting}
#include <concepts>

template<class T>
requires std::integral<T>
T square1(T x) { return x * x; }

template<class T>
T square2(T x) requires std::integral<T> {
    return x * x;
}

template<std::integral T>
T square3(T x) { return x * x; }

std::integral auto square4(std::integral auto x) {
    return x * x;
}
\end{lstlisting}

W ostatnim przykładzie ograniczony jest nie tylko typ argumentu,
ale też \emph{typ wyniku}. Sama dedukcja typu wyniku nadal odbywa
się na podstawie wyrażenia \texttt{x * x}; koncept jest sprawdzany
po jej przeprowadzeniu. Zapis \texttt{std::integral auto} w parametrze
jest skrótem szablonu funkcji, tzw. \emph{abbreviated function template}.

\paragraph{Kiedy którą postać wybrać?}
Dla jednego prostego ograniczenia zapis z ograniczonym
parametrem szablonu, np. \texttt{template<\allowbreak std::integral T>},
jest krótki i czytelny. Jeśli zależność dotyczy kilku parametrów,
najwygodniejsza bywa klauzula \texttt{requires}.
Skrócona postać z \texttt{auto} przydaje się w niewielkich funkcjach
i lambdach.

\section{Dwie role słowa \texttt{requires}}

\begin{center}\small
\begin{tabularx}{\textwidth}{@{}p{0.34\textwidth}X@{}}
\toprule
Zapis & Za co odpowiada \\
\midrule
\code{concept C = ...} & Nadaje nazwę logicznemu ograniczeniu parametrów szablonu. \\
\code{requires C<T>} & Ogranicza udział szablonu w przeciążaniu do typów spełniających \code{C}. \\
\code{requires(T x) { x.f(); }} & Sprawdza poprawność wskazanego wyrażenia dla podstawionego typu. \\
\code{requires requires(...) { ... }} & Pierwsze \code{requires} wprowadza ograniczenie, drugie wyrażenie sprawdzające. \\
\code{static_assert(C<T>)} & Wymusza spełnienie warunku w danym miejscu i zgłasza błąd kompilacji, gdy jest fałszywy. \\
\bottomrule
\end{tabularx}
\end{center}
\label{sec:concepts-dwa-requires}

Identyczny zapis słowny \texttt{requires} występuje w dwóch
różnych konstrukcjach:
\begin{enumerate}
  \item \emph{requires-clause} -- klauzula ograniczająca deklarację;
  \item \emph{requires-expression} -- wyrażenie logiczne, które bada
  poprawność podanych wymagań.
\end{enumerate}

\begin{lstlisting}
#include <concepts>

template<class T>
concept Incrementable = requires(T value) { // requires-expression
    ++value;
};

template<class T>
requires Incrementable<T>                // requires-clause
void increment(T& x) {
    ++x;
}
\end{lstlisting}

Można je połączyć bez nadawania koncepcji osobnej nazwy:

\begin{lstlisting}
template<class T>
requires requires(T x) { ++x; }
void increment_again(T& x) {
    ++x;
}
\end{lstlisting}

Pierwsze \texttt{requires} w ostatnim przykładzie zaczyna klauzulę;
drugie wprowadza wyrażenie sprawdzające wymagania. To poprawna
składnia, choć przy wielokrotnym użyciu lepiej utworzyć nazwaną
koncepcję.

\section{Cztery rodzaje wymagań w \texttt{requires-expression}}

\begin{center}\small
\begin{tabularx}{\textwidth}{@{}p{0.31\textwidth}X@{}}
\toprule
Rodzaj wymagania & Co sprawdza \\
\midrule
Proste: \code{a + b;} & Czy wyrażenie jest poprawne składniowo dla danych typów. \\
Typowe: \code{typename T::value_type;} & Czy wskazany typ zależny istnieje. \\
Złożone: \code{{ f(x) } -> C;} & Czy wyrażenie jest poprawne oraz czy jego typ spełnia ograniczenie \code{C}. \\
Zagnieżdżone: \code{requires C<T>;} & Czy dodatkowy warunek logiczny jest spełniony. \\
\bottomrule
\end{tabularx}
\end{center}
\label{sec:concepts-cztery-wymagania}

Wyrażenie \texttt{requires} umożliwia opisanie czterech rodzajów
warunków. Parametry wpisane w nawiasach \texttt{requires(T x)}
są jedynie \emph{lokalnymi parametrami do sprawdzania wyrażeń};
nie tworzy się dla nich obiektów podczas działania programu.

\subsection{Wymaganie proste (\textit{simple requirement})}

\begin{lstlisting}
#include <ostream>

template<class T>
concept Printable = requires(std::ostream& os, const T& x) {
    os << x;                    // must be well-formed
};
\end{lstlisting}

Wyrażenie \texttt{os << x}
nie jest tutaj wykonywane; kompilator sprawdza jego poprawność.
Proste wymaganie nie nakłada warunku na typ zwracany.

\subsection{Wymaganie typu (\textit{type requirement})}

\begin{lstlisting}
#include <vector>

template<class T>
concept HasValueType = requires {
    typename T::value_type;    // nested type must exist
};

static_assert(HasValueType<std::vector<int>>);
static_assert(!HasValueType<int>);
\end{lstlisting}

Słowo \texttt{typename} jest tu składnikiem składni wymagania,
a nie deklaracją nowego typu.

\subsection{Wymaganie złożone (\textit{compound requirement})}

\begin{lstlisting}
#include <concepts>

template<class T>
concept AddableNoThrow = requires(const T& a, const T& b) {
    { a + b } noexcept -> std::same_as<T>;
};
\end{lstlisting}

Zapis \texttt{\{ wyrazenie \} noexcept -> Koncepcja;}
może sprawdzać jednocześnie: poprawność wyrażenia,
to, że jest ono \texttt{noexcept}, i własność typu
\texttt{decltype((wyrazenie))}. Obie części, \texttt{noexcept}
oraz \texttt{-> Koncepcja}, są opcjonalne.

\begin{lstlisting}
template<class T>
concept PlusReturnsT = requires(const T& a, const T& b) {
    { a + b } -> std::same_as<T>;
};

static_assert(PlusReturnsT<int>);
// A type with a proxy return value may not satisfy this.
\end{lstlisting}

\paragraph{Uwaga na typ i konwersje.}
Koncepcja \texttt{std::same\_as<T>} żąda \emph{dokładnie}
typu \texttt{T}, natomiast koncepcja
\texttt{std::convertible\_\allowbreak to<T>} dopuszcza konwersję.
Dla \texttt{int} operator \texttt{+} daje \texttt{int}, ale np.
dla małych typów całkowitych może zajść promocja do \texttt{int}.

\subsection{Wymaganie zagnieżdżone (\textit{nested requirement})}

\begin{lstlisting}
#include <type_traits>

template<class T>
concept SmallTrivial = requires {
    requires sizeof(T) <= 16;
    requires std::is_trivially_copyable_v<T>;
};

static_assert(SmallTrivial<int>);
\end{lstlisting}

Wymaganie zagnieżdżone zaczyna się kolejnym \texttt{requires}
i musi zawierać warunek stały. Jest wygodne, gdy wymaganie
dotyczy własności typu, a nie określonego wyrażenia.

\subsection{Przykład: połączenie różnych wymagań}

\begin{lstlisting}
#include <concepts>
#include <cstddef>

template<class T>
concept Indexable = requires(T& obj, std::size_t i) {
    typename T::value_type;
    { obj.size() } -> std::convertible_to<std::size_t>;
    obj[i];
};
\end{lstlisting}

To przykład \emph{minimalnego kontraktu składniowego}.
Koncepcja nie dowodzi, że indeks \texttt{i} jest w zakresie,
że operacja jest szybka ani że wynik indeksowania zachowuje się
jak zwykła referencja. Takie własności wymagają dodatkowych założeń.

\section{Standardowe koncepcje z biblioteki C++20}
\label{sec:concepts-standardowe}

Przed napisaniem własnej koncepcji warto sprawdzić, czy istnieje
gotowa w standardowej bibliotece.

\subsection{Typy, konwersje i dziedziczenie}

\begin{lstlisting}
#include <concepts>

static_assert(std::same_as<int, int>);
static_assert(!std::same_as<int, long>);
static_assert(std::integral<unsigned>);
static_assert(std::signed_integral<int>);
static_assert(std::unsigned_integral<unsigned>);
static_assert(std::floating_point<double>);
static_assert(std::convertible_to<int, double>);

struct Base {};
struct Derived : Base {};
static_assert(std::derived_from<Derived, Base>);
\end{lstlisting}

\texttt{std::derived\_from<D,B>} obejmuje publiczne i jednoznaczne
dziedziczenie oraz przypadek tego samego typu. Różni się pod tym
względem od samego \texttt{std::is\_base\_of\_v}.

\subsection{Konstruowalność, kopiowanie, przenoszenie i porównywanie}

Przydatne są m.in. koncepcje wyrażające inicjalizowalność,
konstruowalność, możliwość przenoszenia i kopiowania oraz porównywania:
\begin{lstlisting}
default_initializable    constructible_from
move_constructible       copy_constructible
movable                  copyable
semiregular              regular
equality_comparable      totally_ordered
\end{lstlisting}

\begin{lstlisting}
#include <concepts>
#include <memory>
#include <string>

static_assert(std::regular<int>);
static_assert(std::copyable<std::string>);
static_assert(std::movable<std::unique_ptr<int>>);
static_assert(!std::copyable<std::unique_ptr<int>>);
static_assert(std::totally_ordered<int>);
\end{lstlisting}

Koncepcje te opisują interfejs i zakładają także określone
\emph{prawa semantyczne} (np. sensowne porównywanie). Sam kompilator
nie potrafi udowodnić dowolnej własności semantycznej implementacji.

\subsection{Funkcje wywoływalne i predykaty}

\begin{lstlisting}
#include <concepts>

struct IsPositive {
    bool operator()(int x) const { return x > 0; }
};

static_assert(std::invocable<IsPositive, int>);
static_assert(std::predicate<IsPositive, int>);
static_assert(!std::predicate<IsPositive, void*>);
\end{lstlisting}

\texttt{std::invocable<F,Args...>} sprawdza możliwość wywołania,
natomiast \texttt{std::predicate} wymaga dodatkowo wyniku,
którego można użyć jako warunku logicznego, i wiąże się
z określonymi wymaganiami semantycznymi.

\subsection{Iteratory i zakresy}

Koncepcje iteratorów są udostępniane przez \texttt{<iterator>},
a koncepcje zakresów przez \texttt{<ranges>}.

\begin{lstlisting}
#include <ranges>
#include <vector>
#include <list>

static_assert(std::ranges::random_access_range<
              std::vector<int>>);
static_assert(std::ranges::bidirectional_range<
              std::list<int>>);
static_assert(!std::ranges::random_access_range<
              std::list<int>>);
\end{lstlisting}

W szczególności przydatne są koncepcje \emph{input},
\emph{forward}, \emph{bidirectional}, \emph{random access}
i \emph{contiguous range}. Pozwalają opisywać wymagania
algorytmów bez uzależniania ich od konkretnego kontenera.

\section{Łączenie koncepcji i definiowanie własnych kontraktów}

Wymagania można łączyć operatorami logicznymi
\texttt{\&\&} i \texttt{||}. Często lepiej budować niewielkie,
nazwane koncepcje niż powtarzać jedno długie wyrażenie.

\begin{lstlisting}
#include <concepts>

template<class T>
concept Arithmetic = std::integral<T> || std::floating_point<T>;

template<class T>
concept NonBooleanInteger =
    std::integral<T> && !std::same_as<T, bool>;

template<Arithmetic T>
T cube(T x) { return x * x * x; }

static_assert(Arithmetic<double>);
static_assert(!NonBooleanInteger<bool>);
\end{lstlisting}

Dla połączenia \texttt{A \&\& B} wymaganie \texttt{B} jest
sprawdzane dopiero, gdy \texttt{A} jest spełnione; analogicznie
dla \texttt{A || B} drugi składnik jest potrzebny dopiero,
gdy pierwszy nie jest spełniony. Ta reguła chroni często przed
podstawianiem wyrażeń, które dla danego typu nie mają sensu.

\begin{lstlisting}
#include <concepts>
#include <type_traits>

template<class T>
concept HasIntegralValueType =
    requires { typename T::value_type; } &&
    std::integral<typename T::value_type>;
\end{lstlisting}

W powyższym przykładzie odwołanie do \texttt{T::value\_type}
po prawej stronie jest istotne dopiero po potwierdzeniu,
że typ składowy istnieje. Należy jednak pamiętać, że złożone
reguły podstawiania i normalizacji koncepcji mają swoje szczegóły;
nie każdą źle sformułowaną deklarację naprawi samo zwarcie logiczne.

\section{Koncepcje z wieloma parametrami}

Koncepcja nie musi opisywać własności jednego typu.
Może także opisywać \emph{relację między typami}.

\begin{lstlisting}
#include <concepts>

template<class From, class To>
concept SameType = std::same_as<From, To>;

template<class Derived, class Base>
concept PublicChildOf = std::derived_from<Derived, Base>;

template<class T, class U>
concept AddableTogether = requires(const T& a, const U& b) {
    a + b;
};

static_assert(AddableTogether<int, double>);
\end{lstlisting}

Wymienione koncepcje \texttt{SameType} i
\texttt{PublicChildOf} służą tu jedynie ilustracji składni.
W praktyce używaj bezpośrednio standardowych
\texttt{std::same\_as} i \texttt{std::derived\_from},
jeżeli nie dodajesz do nich istotnego znaczenia domenowego.

\paragraph{Wstawianie argumentu przy ograniczonym parametrze.}

\begin{lstlisting}
struct Animal {};

template<std::derived_from<Animal> T>
void inspect(const T&) {}

// Equivalent constraint: std::derived_from<T, Animal>
\end{lstlisting}

W zapisie powyżej typ \texttt{T} jest podstawiany
\emph{na pierwszej pozycji} koncepcji, zatem
chodzi o warunek z poniższego przykładu:
\begin{lstlisting}
std::derived_from<T, Animal>
\end{lstlisting}

\section{Przeciążanie funkcji i wybór bardziej szczegółowego ograniczenia}
\label{sec:concepts-przeciazanie}

Ograniczenia uczestniczą w wyborze przeciążenia. Jeśli dwie funkcje
są dopasowane do argumentów, kompilator może preferować tę,
której ograniczenie jest bardziej szczegółowe (ang.
\emph{subsumption}).

\begin{lstlisting}
#include <concepts>
#include <string_view>

template<class T>
requires std::integral<T>
constexpr std::string_view category(T) {
    return "integral";
}

template<class T>
requires std::signed_integral<T>
constexpr std::string_view category(T) {
    return "signed integral";
}

static_assert(category(12) == "signed integral");
static_assert(category(12u) == "integral");
\end{lstlisting}

\texttt{std::signed\_integral} jest zbudowane na
\texttt{std::integral} i dodaje dodatkowe wymaganie. Dzięki temu
bardziej szczegółowe przeciążenie zostaje wybrane dla \texttt{int}.

\subsection{Ważna pułapka: podobny zapis to nie to samo ograniczenie}

Porządek ograniczeń nie jest ogólnym dowodzeniem implikacji
matematycznych. Kompilator porównuje znormalizowane ograniczenia
atomowe i ich \emph{tożsamość}, a nie tylko sens pozornie podobnych
wyrażeń.

\begin{lstlisting}
#include <type_traits>

template<class T>
concept HasSize = requires(const T& x) { x.size(); };

template<class T>
concept HasSizeAndCopy =
    HasSize<T> && std::is_copy_constructible_v<T>;

template<class T> requires HasSize<T>
int classify(const T&) { return 1; }

template<class T> requires HasSizeAndCopy<T>
int classify(const T&) { return 2; }
\end{lstlisting}

Druga koncepcja \emph{wykorzystuje tę samą} koncepcję
\texttt{HasSize}, dlatego kompilator rozpoznaje wspólny
składnik. Jeśli powtórzymy ręcznie nawet identycznie wyglądające
wymaganie opisujące metodę \texttt{size()},
nie mamy automatycznej gwarancji identyczności ograniczeń
atomowych.

\paragraph{Rekomendacja.}
Twórz proste nazwane koncepcje i buduj bardziej szczegółowe
z ich wykorzystaniem. Ułatwia to konserwację kodu i
przewidywanie wyboru przeciążeń.

\section{Koncepcje dla szablonów klas i ich specjalizacji}

Ograniczać można nie tylko funkcje, lecz również szablony klas.

\begin{lstlisting}
#include <concepts>

template<std::integral T>
class Counter {
    T value_{};
public:
    void add(T delta) { value_ += delta; }
    T value() const { return value_; }
};

// Counter<int> a;       // OK
// Counter<double> b;    // ERROR: constraint not satisfied
\end{lstlisting}

Ograniczenia mogą też sterować wyborem częściowej specjalizacji.

\begin{lstlisting}
#include <concepts>
#include <string_view>

template<class T>
struct TypeLabel {
    static constexpr std::string_view value = "other";
};

template<std::integral T>
struct TypeLabel<T> {
    static constexpr std::string_view value = "integer";
};

template<std::floating_point T>
struct TypeLabel<T> {
    static constexpr std::string_view value = "floating";
};

static_assert(TypeLabel<int>::value == "integer");
static_assert(TypeLabel<double>::value == "floating");
\end{lstlisting}

Nie używaj jednak złożonego systemu specjalizacji, jeżeli
prosta funkcja przeciążona lub \texttt{if constexpr}
lepiej wyraża zamysł.

\section{Koncepcje i pakiety parametrów: kontynuacja rozdziału 4}
\label{sec:concepts-variadic}

Koncepcje bardzo dobrze współpracują z \emph{variadic templates}.
Możemy nakładać warunek na każdy element pakietu typów
za pomocą \emph{fold expression}.

\begin{lstlisting}
#include <concepts>

template<class... Ts>
requires (std::integral<Ts> && ...)
constexpr auto sum_integers(Ts... xs) {
    return (0 + ... + xs);
}

static_assert(sum_integers() == 0);
static_assert(sum_integers(1, 2, 3) == 6);
// sum_integers(1, 2.5);    // ERROR: double not integral
\end{lstlisting}

Zapis \texttt{(P<Ts> \&\& ...)} reprezentuje
koniunkcję wymagań dla wszystkich elementów pakietu.
Dla pustego pakietu taki fold logiczny daje \texttt{true};
nie oznacza to jednak, że każde ciało funkcji poradzi sobie
z zerową liczbą argumentów. W przykładzie użyliśmy
\texttt{0} jako elementu neutralnego sumowania.

\subsection{Ograniczony pakiet typów i uniwersalne przekazywanie}

\begin{lstlisting}
#include <concepts>
#include <utility>

template<std::integral... Ts>
void accept_integers(Ts&&... xs) {
    ((void)std::forward<Ts>(xs), ...); // illustrative use
}
\end{lstlisting}

Zapis \texttt{std::integral... Ts} również ogranicza wszystkie
elementy pakietu. Uwaga: przy dedukcji z przekazywanymi referencjami
\texttt{Ts\&\&}, typ \texttt{Ts} może być \emph{referencją},
gdy argument jest lvalue. W takim przypadku
\texttt{std::integral<Ts>} może być fałszywe.
Jeżeli chcemy akceptować obie kategorie wartości, często
właściwsze jest użycie \texttt{std::remove\_cvref\_t}:

\begin{lstlisting}
#include <concepts>
#include <type_traits>
#include <utility>

template<class... Ts>
requires (std::integral<std::remove_cvref_t<Ts>> && ...)
void accept_any_int(Ts&&... xs) {
    ((void)xs, ...); // demonstration; no ownership transfer
}
\end{lstlisting}

Ta różnica jest szczególnie ważna przy projektowaniu generycznych
wrapperów, które mają przyjmować także obiekty \texttt{const}
i referencje.

\section{Concepts a \texttt{if constexpr}, \texttt{static\_assert} i SFINAE}
\label{sec:concepts-porownanie}

Są to narzędzia pokrewne, ale mają różne cele.

\subsection{\texttt{if constexpr}: wybór implementacji}

\begin{lstlisting}
#include <concepts>
#include <string>
#include <string_view>

template<class T>
std::string describe(const T& value) {
    if constexpr (std::integral<T>) {
        return "an integer";
    } else if constexpr (std::floating_point<T>) {
        return "a floating-point value";
    } else {
        return "another type";
    }
}
\end{lstlisting}

\texttt{if constexpr} umożliwia wybranie gałęzi kodu
\emph{wewnątrz} szablonu. Nie tworzy samo przez się
ograniczonego interfejsu. Z kolei \texttt{requires}
może odrzucić niewłaściwe wywołanie jeszcze na etapie
dobierania przeciążeń.

\subsection{\texttt{static\_assert}: własny komunikat}

\begin{lstlisting}
#include <concepts>

template<class T>
void save_number(const T& value) {
    static_assert(std::integral<T>,
                  "save_number expects an integral type");
    // ...
}
\end{lstlisting}

\texttt{static\_assert} jest użyteczne, gdy potrzebujemy
jednoznacznego komunikatu albo sprawdzamy niezmiennik wewnątrz
implementacji. Nie zastępuje jednak ograniczenia uczestniczącego
w rozstrzyganiu przeciążeń.

\subsection{Przed C++20: \texttt{std::enable\_if} i idiom wykrywania}

\begin{lstlisting}
#include <type_traits>

// C++11/14/17 technique:
template<class T,
         std::enable_if_t<std::is_integral_v<T>, int> = 0>
T old_square(T x) {
    return x * x;
}

// C++20:
#include <concepts>
template<std::integral T>
T new_square(T x) {
    return x * x;
}
\end{lstlisting}

Dla nowego kodu C++20 najczęściej korzystniej wybrać
koncepcję: intencja staje się czytelna i spójna z mechanizmem
porządkowania przeciążeń. \texttt{std::enable\_if}
nadal ma znaczenie w bibliotekach wspierających starsze
wersje standardu.

\section{Praktyczny przykład: średnia z zakresu liczb}
\label{sec:concepts-praktyka}

Zamiast pisać funkcję tylko dla \texttt{std::vector<double>},
możemy określić dokładne wymagania wobec \emph{zakresu}.

\begin{lstlisting}
#include <concepts>
#include <cstddef>
#include <optional>
#include <ranges>
#include <vector>
#include <array>

template<class R>
concept NumericInputRange =
    std::ranges::input_range<R> &&
    std::convertible_to<std::ranges::range_reference_t<R>, double>;

template<NumericInputRange R>
std::optional<double> average(R&& range) {
    double sum = 0.0;
    std::size_t count = 0;
    for (auto&& item : range) {
        sum += static_cast<double>(item);
        ++count;
    }
    if (count == 0) return std::nullopt;
    return sum / static_cast<double>(count);
}

int main() {
    std::vector<int> v{2, 4, 6};
    std::array<double, 2> a{3.0, 7.0};
    auto x = average(v); // 4.0
    auto y = average(a); // 5.0
    auto z = average(std::vector<int>{}); // empty optional
    return !(x && y && !z && *x == 4.0 && *y == 5.0);
}
\end{lstlisting}

Funkcja działa dla różnych kontenerów i zakresów dostarczających
iterację jednokierunkową. Warunek \texttt{input\_range}
wyraża strukturę danych, a \texttt{convertible\_to}
-- oczekiwany rodzaj elementów. Zastosowaliśmy
\texttt{std::optional<double>}, aby pusty zakres nie powodował
dzielenia przez zero. Dla zastosowań numerycznych należy
dodatkowo rozważyć stabilność sumowania, przepełnienia i
znaczenie konwersji do \texttt{double}.

\paragraph{Kiedy \texttt{ranges}?}
Koncepcje standardowych zakresów są szczególnie użyteczne,
gdy funkcja operuje na sekwencji, ale nie wymaga konkretnie
\texttt{std::vector} czy \texttt{std::list}. Zamiast wymuszać
silniejsze wymaganie \texttt{random\_access\_range},
wybieraj najsłabszą koncepcję, która wystarcza do działania
algorytmu.

\section{Koncepcje dla lambd i algorytmów}

Funkcje anonimowe również mogą mieć ograniczone parametry.

\begin{lstlisting}
#include <concepts>

int main() {
    auto twice = [](std::integral auto x) { return x + x; };
    return twice(5) != 10;
}
\end{lstlisting}

Można też tworzyć generyczne algorytmy przyjmujące funkcje
spełniające \texttt{std::predicate}.

\begin{lstlisting}
#include <concepts>
#include <cstddef>
#include <vector>

template<class T, class Pred>
requires std::predicate<Pred&, const T&>
std::size_t count_if_custom(const std::vector<T>& values,
                            Pred pred) {
    std::size_t count = 0;
    for (const T& x : values)
        if (pred(x)) ++count;
    return count;
}
\end{lstlisting}

Zwróć uwagę na argumenty typu w
\texttt{std::predicate<Pred\&, const T\&>}:
sprawdzają możliwość wywołania \emph{konkretnego rodzaju}
predykatu z \emph{konkretnym rodzajem} argumentu.
Kategoria wartości ma znaczenie i nie należy bezrefleksyjnie
wstawiać tu samych nazw typów.

\section{Granice gwarancji: najczęstsze pułapki}
\label{sec:concepts-pulapki}

\subsection{Concept nie sprawdza dowolnej własności semantycznej}

\begin{lstlisting}
#include <concepts>

struct Strange {
    int value{};
    bool operator==(const Strange&) const { return false; }
};

static_assert(std::equality_comparable<Strange>);
// Syntactically valid, but equality lacks reflexivity.
\end{lstlisting}

Kompilator akceptuje składnię porównania, lecz
\texttt{a == a} jest zawsze fałszywe. Kod nie spełnia
matematycznych oczekiwań związanych z relacją równości.
To obowiązek autora typu, niezależny od samego mechanizmu
\texttt{concepts}.

\subsection{Koncepcje nie zamieniają błędów w poprawny kod}

Wyrażenia \texttt{requires} pozwalają w kontekście zależnym
od parametru szablonu traktować niepoprawne podstawienia
jako niespełnione wymagania. Nie oznacza to jednak, że
\emph{każdy} błędny fragment C++ staje się automatycznie
\texttt{false}. Niezależne od parametru, trwale niepoprawne
konstrukcje mogą prowadzić do błędu kompilacji.

\subsection{Różnica między \texttt{T}, \texttt{T\&} i \texttt{const T\&}}

\begin{lstlisting}
#include <concepts>
#include <type_traits>

template<class T>
concept PlainInteger =
    std::integral<std::remove_cvref_t<T>>;

static_assert(!std::integral<int&>);
static_assert(PlainInteger<int&>);
static_assert(PlainInteger<const int&>);
\end{lstlisting}

Przy dedukcji typu generycznego parametru
\texttt{T\&\&} ta różnica często decyduje o tym,
czy z pozoru poprawne wywołanie spełni koncepcję.

\subsection{\texttt{std::same\_as} czy \texttt{std::convertible\_to}?}

\begin{lstlisting}
#include <concepts>

template<class T>
concept StrictAdd = requires(T a, T b) {
    { a + b } -> std::same_as<T>;
};

template<class T>
concept FlexibleAdd = requires(T a, T b) {
    { a + b } -> std::convertible_to<T>;
};

static_assert(!StrictAdd<short>);  // short + short -> int
static_assert(FlexibleAdd<short>); // int -> short allowed
\end{lstlisting}

Koncepcja \texttt{FlexibleAdd} nie oznacza \emph{bezstratnej}
konwersji; \texttt{int} można niejawnie przekształcić do
\texttt{short}, tracąc informację. Nazwa koncepcji powinna
odzwierciedlać rzeczywistą, a nie postulowaną gwarancję.

\subsection{Funkcja ograniczona nie zawsze jest optymalniejsza}

Koncepcje opisują poprawność i wybór instancji szablonów.
Nie wymuszają \emph{inline}, wektoryzacji, eliminacji kopii
ani lepszego kodu maszynowego. Można je stosować bez obietnicy
przyspieszenia wykonywania programu.

\section{Zasady projektowania: kiedy używać, a kiedy nie?}

\begin{center}
\small
\begin{tabular}{|p{0.34\textwidth}|p{0.55\textwidth}|}
\hline
\textbf{Sytuacja} & \textbf{Zalecane rozwiązanie} \\
\hline
Funkcja generyczna wymaga konkretnej kategorii typów & Nazwany \texttt{concept}, np. \texttt{std::integral} \\
\hline
Wymagania dotyczą istnienia operacji lub typu składowego & Własny \texttt{requires-expression} \\
\hline
Ważne są różne wersje algorytmu dla różnych typów & Ograniczone przeciążenia lub specjalizacje \\
\hline
Jedna implementacja ma kilka gałęzi zależnych od typu & \texttt{if constexpr} \\
\hline
Sprawdzasz założenie wewnątrz ciała funkcji & \texttt{static\_assert} z komunikatem \\
\hline
Weryfikujesz jednakowy warunek dla całego pakietu & Fold \texttt{(Concept<Ts> \&\& ...)} \\
\hline
Potrzebujesz iteracji po sekwencji & \texttt{std::ranges::input\_range} lub silniejszy koncept \\
\hline
Funkcja działa dla jednego, konkretnego typu & Zwykła funkcja często jest prostsza niż szablon \\
\hline
Projekt musi działać także w C++17 & \texttt{type\_traits} i (w razie potrzeby) SFINAE \\
\hline
Wymagasz warunku wykonywania programu, np. \texttt{n>0} & Zwykłe sprawdzenie w czasie działania; koncept go nie zastępuje \\
\hline
\end{tabular}
\end{center}

\section{Podsumowanie i pytania kontrolne}

\begin{enumerate}
\item \texttt{concept} nazywa predykat nad parametrami szablonu,
sprawdzany w czasie kompilacji.
\item \texttt{requires-clause} ogranicza deklarację;
\texttt{requires-expression} jest wyrażeniem badającym wymagania.
\item Istnieją cztery formy wymagań: proste, typu, złożone
i zagnieżdżone.
\item \texttt{std::same\_as}, \texttt{std::convertible\_to},
\texttt{std::integral}, \texttt{std::predicate} i koncepcje zakresów
powinny być zwykle pierwszym wyborem.
\item Ograniczenia mogą wpływać na wybór przeciążenia, ale mechanizm
porządkowania nie jest dowolnym dowodzeniem implikacji logicznej.
\item Nazwane koncepcje warto wykorzystywać ponownie przy
budowaniu koncepcji bardziej szczegółowych.
\item \texttt{std::remove\_cvref\_t<T>} ma często znaczenie
przy koncepcjach typów dedukowanych z referencji przekazujących.
\item Koncepcje opisują przede wszystkim wymagania składniowe,
a dodatkowe prawa semantyczne muszą respektować autorzy typów.
\end{enumerate}

\paragraph{Pytania do samodzielnego sprawdzenia.}
Zastanów się, co kompilator zaakceptuje i dlaczego:

\begin{lstlisting}
#include <concepts>
#include <type_traits>

static_assert(std::integral<bool>);
static_assert(!std::integral<int&>);
static_assert(std::integral<std::remove_cvref_t<const int&>>);
static_assert(!std::same_as<int, long>);

// 1. Why does StrictAdd<short> fail?
// 2. Why can a concept be true for a logically broken ==?
// 3. How do requires-clause and requires-expression differ?
// 4. Which overload of category(5) is chosen?
// 5. Why is NumericInputRange more general than vector<double>?
\end{lstlisting}

\paragraph{Źródła i dalsza lektura.}
Oficjalne reguły języka określa standard ISO C++, a przystępne
opisy i przykłady znajdują się w serwisie \emph{cppreference}.
Przydatne adresy:
\begin{lstlisting}
en.cppreference.com/cpp/language/constraints
en.cppreference.com/cpp/language/requires
en.cppreference.com/cpp/concepts
\end{lstlisting}


\chapter[Concepts i requires w C++]{Koncepcje i ograniczenia szablonów (\textit{concepts} i \texttt{requires})}
\label{chap:concepts-cpp}

\section{Po co wprowadzono koncepcje?}

Szablony C++ pozwalają pisać kod niezależny od konkretnego typu.
Tradycyjny szablon deklaruje jednak zwykle \emph{co kompilator ma zrobić},
nie wyjaśniając wprost, \emph{jakich własności} wymaga od argumentów.
Jeśli operacja jest niepoprawna, błąd może zostać wykryty dopiero w ciele
szablonu, a komunikat diagnostyczny bywa długi.

\begin{lstlisting}
template<class T>
T twice(const T& x) {
    return x + x;           // requires operator+
}

// twice(7);                // OK
// twice(SomeType{});       // may fail inside function body
\end{lstlisting}

Od C++20 można umieścić wymagania na poziomie interfejsu:

\begin{lstlisting}
#include <concepts>

template<class T>
concept Addable = requires(const T& a, const T& b) {
    { a + b } -> std::convertible_to<T>;
};

template<Addable T>
T twice(const T& x) {
    return x + x;
}
\end{lstlisting}

Koncepcja (\textit{concept}) to nazwany predykat sprawdzany podczas
kompilacji, opisujący wymagania wobec argumentów szablonu.
Może dotyczyć istnienia operatorów, typów składowych, własności typów,
kategorii iteratorów i wielu innych cech. \emph{Constraints}
(ograniczenia) wpływają zarówno na dopuszczalność konkretyzacji,
jak i na wybór przeciążenia funkcji.

\paragraph{Wersja języka.}
W całym rozdziale zakładamy \textbf{C++20} oraz kompilację
np. poleceniem \texttt{g++ -std=c++20 -Wall -Wextra -Wpedantic main.cpp}.
Definicje koncepcji standardowych znajdują się przede wszystkim
w nagłówkach \texttt{<concepts>}, \texttt{<iterator>} i
\texttt{<ranges>}.

\section{Pierwsza koncepcja: składnia i zastosowanie}
\label{sec:concepts-pierwsza}

Definicja nazwanej koncepcji ma postać
\texttt{template<...> concept Nazwa = warunek;}. Warunek powinien
być wyrażeniem stałym typu \texttt{bool}. Definicja koncepcji
występuje na poziomie przestrzeni nazw; nie jest zwykłą funkcją.

\begin{lstlisting}
#include <concepts>
#include <type_traits>

template<class T>
concept Integer = std::is_integral_v<T>;

static_assert(Integer<int>);
static_assert(Integer<bool>);    // bool is integral
static_assert(!Integer<double>);
\end{lstlisting}

Koncepcję można sprawdzić wyrażeniem \texttt{Integer<int>} lub
\texttt{static\_assert}, ale nie należy jej mylić ze zwykłą
zmienną. W szczególności \texttt{concept} nie daje się jawnie
specjalizować w sposób analogiczny do szablonu klasy.

\subsection{Cztery podstawowe sposoby ograniczania funkcji}

Poniższe deklaracje opisują ten sam warunek: funkcja jest przeznaczona
dla typów całkowitych w sensie \texttt{std::integral}.

\begin{lstlisting}
#include <concepts>

template<class T>
requires std::integral<T>
T square1(T x) { return x * x; }

template<class T>
T square2(T x) requires std::integral<T> {
    return x * x;
}

template<std::integral T>
T square3(T x) { return x * x; }

std::integral auto square4(std::integral auto x) {
    return x * x;
}
\end{lstlisting}

W ostatnim przykładzie ograniczony jest nie tylko typ argumentu,
ale też \emph{typ wyniku}. Sama dedukcja typu wyniku nadal odbywa
się na podstawie wyrażenia \texttt{x * x}; koncept jest sprawdzany
po jej przeprowadzeniu. Zapis \texttt{std::integral auto} w parametrze
jest skrótem szablonu funkcji, tzw. \emph{abbreviated function template}.

\paragraph{Kiedy którą postać wybrać?}
Dla jednego prostego ograniczenia zapis z ograniczonym
parametrem szablonu, np. \texttt{template<\allowbreak std::integral T>},
jest krótki i czytelny. Jeśli zależność dotyczy kilku parametrów,
najwygodniejsza bywa klauzula \texttt{requires}.
Skrócona postać z \texttt{auto} przydaje się w niewielkich funkcjach
i lambdach.

\section{Dwie role słowa \texttt{requires}}

\begin{center}\small
\begin{tabularx}{\textwidth}{@{}p{0.34\textwidth}X@{}}
\toprule
Zapis & Za co odpowiada \\
\midrule
\code{concept C = ...} & Nadaje nazwę logicznemu ograniczeniu parametrów szablonu. \\
\code{requires C<T>} & Ogranicza udział szablonu w przeciążaniu do typów spełniających \code{C}. \\
\code{requires(T x) { x.f(); }} & Sprawdza poprawność wskazanego wyrażenia dla podstawionego typu. \\
\code{requires requires(...) { ... }} & Pierwsze \code{requires} wprowadza ograniczenie, drugie wyrażenie sprawdzające. \\
\code{static_assert(C<T>)} & Wymusza spełnienie warunku w danym miejscu i zgłasza błąd kompilacji, gdy jest fałszywy. \\
\bottomrule
\end{tabularx}
\end{center}
\label{sec:concepts-dwa-requires}

Identyczny zapis słowny \texttt{requires} występuje w dwóch
różnych konstrukcjach:
\begin{enumerate}
  \item \emph{requires-clause} -- klauzula ograniczająca deklarację;
  \item \emph{requires-expression} -- wyrażenie logiczne, które bada
  poprawność podanych wymagań.
\end{enumerate}

\begin{lstlisting}
#include <concepts>

template<class T>
concept Incrementable = requires(T value) { // requires-expression
    ++value;
};

template<class T>
requires Incrementable<T>                // requires-clause
void increment(T& x) {
    ++x;
}
\end{lstlisting}

Można je połączyć bez nadawania koncepcji osobnej nazwy:

\begin{lstlisting}
template<class T>
requires requires(T x) { ++x; }
void increment_again(T& x) {
    ++x;
}
\end{lstlisting}

Pierwsze \texttt{requires} w ostatnim przykładzie zaczyna klauzulę;
drugie wprowadza wyrażenie sprawdzające wymagania. To poprawna
składnia, choć przy wielokrotnym użyciu lepiej utworzyć nazwaną
koncepcję.

\section{Cztery rodzaje wymagań w \texttt{requires-expression}}

\begin{center}\small
\begin{tabularx}{\textwidth}{@{}p{0.31\textwidth}X@{}}
\toprule
Rodzaj wymagania & Co sprawdza \\
\midrule
Proste: \code{a + b;} & Czy wyrażenie jest poprawne składniowo dla danych typów. \\
Typowe: \code{typename T::value_type;} & Czy wskazany typ zależny istnieje. \\
Złożone: \code{{ f(x) } -> C;} & Czy wyrażenie jest poprawne oraz czy jego typ spełnia ograniczenie \code{C}. \\
Zagnieżdżone: \code{requires C<T>;} & Czy dodatkowy warunek logiczny jest spełniony. \\
\bottomrule
\end{tabularx}
\end{center}
\label{sec:concepts-cztery-wymagania}

Wyrażenie \texttt{requires} umożliwia opisanie czterech rodzajów
warunków. Parametry wpisane w nawiasach \texttt{requires(T x)}
są jedynie \emph{lokalnymi parametrami do sprawdzania wyrażeń};
nie tworzy się dla nich obiektów podczas działania programu.

\subsection{Wymaganie proste (\textit{simple requirement})}

\begin{lstlisting}
#include <ostream>

template<class T>
concept Printable = requires(std::ostream& os, const T& x) {
    os << x;                    // must be well-formed
};
\end{lstlisting}

Wyrażenie \texttt{os << x}
nie jest tutaj wykonywane; kompilator sprawdza jego poprawność.
Proste wymaganie nie nakłada warunku na typ zwracany.

\subsection{Wymaganie typu (\textit{type requirement})}

\begin{lstlisting}
#include <vector>

template<class T>
concept HasValueType = requires {
    typename T::value_type;    // nested type must exist
};

static_assert(HasValueType<std::vector<int>>);
static_assert(!HasValueType<int>);
\end{lstlisting}

Słowo \texttt{typename} jest tu składnikiem składni wymagania,
a nie deklaracją nowego typu.

\subsection{Wymaganie złożone (\textit{compound requirement})}

\begin{lstlisting}
#include <concepts>

template<class T>
concept AddableNoThrow = requires(const T& a, const T& b) {
    { a + b } noexcept -> std::same_as<T>;
};
\end{lstlisting}

Zapis \texttt{\{ wyrazenie \} noexcept -> Koncepcja;}
może sprawdzać jednocześnie: poprawność wyrażenia,
to, że jest ono \texttt{noexcept}, i własność typu
\texttt{decltype((wyrazenie))}. Obie części, \texttt{noexcept}
oraz \texttt{-> Koncepcja}, są opcjonalne.

\begin{lstlisting}
template<class T>
concept PlusReturnsT = requires(const T& a, const T& b) {
    { a + b } -> std::same_as<T>;
};

static_assert(PlusReturnsT<int>);
// A type with a proxy return value may not satisfy this.
\end{lstlisting}

\paragraph{Uwaga na typ i konwersje.}
Koncepcja \texttt{std::same\_as<T>} żąda \emph{dokładnie}
typu \texttt{T}, natomiast koncepcja
\texttt{std::convertible\_\allowbreak to<T>} dopuszcza konwersję.
Dla \texttt{int} operator \texttt{+} daje \texttt{int}, ale np.
dla małych typów całkowitych może zajść promocja do \texttt{int}.

\subsection{Wymaganie zagnieżdżone (\textit{nested requirement})}

\begin{lstlisting}
#include <type_traits>

template<class T>
concept SmallTrivial = requires {
    requires sizeof(T) <= 16;
    requires std::is_trivially_copyable_v<T>;
};

static_assert(SmallTrivial<int>);
\end{lstlisting}

Wymaganie zagnieżdżone zaczyna się kolejnym \texttt{requires}
i musi zawierać warunek stały. Jest wygodne, gdy wymaganie
dotyczy własności typu, a nie określonego wyrażenia.

\subsection{Przykład: połączenie różnych wymagań}

\begin{lstlisting}
#include <concepts>
#include <cstddef>

template<class T>
concept Indexable = requires(T& obj, std::size_t i) {
    typename T::value_type;
    { obj.size() } -> std::convertible_to<std::size_t>;
    obj[i];
};
\end{lstlisting}

To przykład \emph{minimalnego kontraktu składniowego}.
Koncepcja nie dowodzi, że indeks \texttt{i} jest w zakresie,
że operacja jest szybka ani że wynik indeksowania zachowuje się
jak zwykła referencja. Takie własności wymagają dodatkowych założeń.

\section{Standardowe koncepcje z biblioteki C++20}
\label{sec:concepts-standardowe}

Przed napisaniem własnej koncepcji warto sprawdzić, czy istnieje
gotowa w standardowej bibliotece.

\subsection{Typy, konwersje i dziedziczenie}

\begin{lstlisting}
#include <concepts>

static_assert(std::same_as<int, int>);
static_assert(!std::same_as<int, long>);
static_assert(std::integral<unsigned>);
static_assert(std::signed_integral<int>);
static_assert(std::unsigned_integral<unsigned>);
static_assert(std::floating_point<double>);
static_assert(std::convertible_to<int, double>);

struct Base {};
struct Derived : Base {};
static_assert(std::derived_from<Derived, Base>);
\end{lstlisting}

\texttt{std::derived\_from<D,B>} obejmuje publiczne i jednoznaczne
dziedziczenie oraz przypadek tego samego typu. Różni się pod tym
względem od samego \texttt{std::is\_base\_of\_v}.

\subsection{Konstruowalność, kopiowanie, przenoszenie i porównywanie}

Przydatne są m.in. koncepcje wyrażające inicjalizowalność,
konstruowalność, możliwość przenoszenia i kopiowania oraz porównywania:
\begin{lstlisting}
default_initializable    constructible_from
move_constructible       copy_constructible
movable                  copyable
semiregular              regular
equality_comparable      totally_ordered
\end{lstlisting}

\begin{lstlisting}
#include <concepts>
#include <memory>
#include <string>

static_assert(std::regular<int>);
static_assert(std::copyable<std::string>);
static_assert(std::movable<std::unique_ptr<int>>);
static_assert(!std::copyable<std::unique_ptr<int>>);
static_assert(std::totally_ordered<int>);
\end{lstlisting}

Koncepcje te opisują interfejs i zakładają także określone
\emph{prawa semantyczne} (np. sensowne porównywanie). Sam kompilator
nie potrafi udowodnić dowolnej własności semantycznej implementacji.

\subsection{Funkcje wywoływalne i predykaty}

\begin{lstlisting}
#include <concepts>

struct IsPositive {
    bool operator()(int x) const { return x > 0; }
};

static_assert(std::invocable<IsPositive, int>);
static_assert(std::predicate<IsPositive, int>);
static_assert(!std::predicate<IsPositive, void*>);
\end{lstlisting}

\texttt{std::invocable<F,Args...>} sprawdza możliwość wywołania,
natomiast \texttt{std::predicate} wymaga dodatkowo wyniku,
którego można użyć jako warunku logicznego, i wiąże się
z określonymi wymaganiami semantycznymi.

\subsection{Iteratory i zakresy}

Koncepcje iteratorów są udostępniane przez \texttt{<iterator>},
a koncepcje zakresów przez \texttt{<ranges>}.

\begin{lstlisting}
#include <ranges>
#include <vector>
#include <list>

static_assert(std::ranges::random_access_range<
              std::vector<int>>);
static_assert(std::ranges::bidirectional_range<
              std::list<int>>);
static_assert(!std::ranges::random_access_range<
              std::list<int>>);
\end{lstlisting}

W szczególności przydatne są koncepcje \emph{input},
\emph{forward}, \emph{bidirectional}, \emph{random access}
i \emph{contiguous range}. Pozwalają opisywać wymagania
algorytmów bez uzależniania ich od konkretnego kontenera.

\section{Łączenie koncepcji i definiowanie własnych kontraktów}

Wymagania można łączyć operatorami logicznymi
\texttt{\&\&} i \texttt{||}. Często lepiej budować niewielkie,
nazwane koncepcje niż powtarzać jedno długie wyrażenie.

\begin{lstlisting}
#include <concepts>

template<class T>
concept Arithmetic = std::integral<T> || std::floating_point<T>;

template<class T>
concept NonBooleanInteger =
    std::integral<T> && !std::same_as<T, bool>;

template<Arithmetic T>
T cube(T x) { return x * x * x; }

static_assert(Arithmetic<double>);
static_assert(!NonBooleanInteger<bool>);
\end{lstlisting}

Dla połączenia \texttt{A \&\& B} wymaganie \texttt{B} jest
sprawdzane dopiero, gdy \texttt{A} jest spełnione; analogicznie
dla \texttt{A || B} drugi składnik jest potrzebny dopiero,
gdy pierwszy nie jest spełniony. Ta reguła chroni często przed
podstawianiem wyrażeń, które dla danego typu nie mają sensu.

\begin{lstlisting}
#include <concepts>
#include <type_traits>

template<class T>
concept HasIntegralValueType =
    requires { typename T::value_type; } &&
    std::integral<typename T::value_type>;
\end{lstlisting}

W powyższym przykładzie odwołanie do \texttt{T::value\_type}
po prawej stronie jest istotne dopiero po potwierdzeniu,
że typ składowy istnieje. Należy jednak pamiętać, że złożone
reguły podstawiania i normalizacji koncepcji mają swoje szczegóły;
nie każdą źle sformułowaną deklarację naprawi samo zwarcie logiczne.

\section{Koncepcje z wieloma parametrami}

Koncepcja nie musi opisywać własności jednego typu.
Może także opisywać \emph{relację między typami}.

\begin{lstlisting}
#include <concepts>

template<class From, class To>
concept SameType = std::same_as<From, To>;

template<class Derived, class Base>
concept PublicChildOf = std::derived_from<Derived, Base>;

template<class T, class U>
concept AddableTogether = requires(const T& a, const U& b) {
    a + b;
};

static_assert(AddableTogether<int, double>);
\end{lstlisting}

Wymienione koncepcje \texttt{SameType} i
\texttt{PublicChildOf} służą tu jedynie ilustracji składni.
W praktyce używaj bezpośrednio standardowych
\texttt{std::same\_as} i \texttt{std::derived\_from},
jeżeli nie dodajesz do nich istotnego znaczenia domenowego.

\paragraph{Wstawianie argumentu przy ograniczonym parametrze.}

\begin{lstlisting}
struct Animal {};

template<std::derived_from<Animal> T>
void inspect(const T&) {}

// Equivalent constraint: std::derived_from<T, Animal>
\end{lstlisting}

W zapisie powyżej typ \texttt{T} jest podstawiany
\emph{na pierwszej pozycji} koncepcji, zatem
chodzi o warunek z poniższego przykładu:
\begin{lstlisting}
std::derived_from<T, Animal>
\end{lstlisting}

\section{Przeciążanie funkcji i wybór bardziej szczegółowego ograniczenia}
\label{sec:concepts-przeciazanie}

Ograniczenia uczestniczą w wyborze przeciążenia. Jeśli dwie funkcje
są dopasowane do argumentów, kompilator może preferować tę,
której ograniczenie jest bardziej szczegółowe (ang.
\emph{subsumption}).

\begin{lstlisting}
#include <concepts>
#include <string_view>

template<class T>
requires std::integral<T>
constexpr std::string_view category(T) {
    return "integral";
}

template<class T>
requires std::signed_integral<T>
constexpr std::string_view category(T) {
    return "signed integral";
}

static_assert(category(12) == "signed integral");
static_assert(category(12u) == "integral");
\end{lstlisting}

\texttt{std::signed\_integral} jest zbudowane na
\texttt{std::integral} i dodaje dodatkowe wymaganie. Dzięki temu
bardziej szczegółowe przeciążenie zostaje wybrane dla \texttt{int}.

\subsection{Ważna pułapka: podobny zapis to nie to samo ograniczenie}

Porządek ograniczeń nie jest ogólnym dowodzeniem implikacji
matematycznych. Kompilator porównuje znormalizowane ograniczenia
atomowe i ich \emph{tożsamość}, a nie tylko sens pozornie podobnych
wyrażeń.

\begin{lstlisting}
#include <type_traits>

template<class T>
concept HasSize = requires(const T& x) { x.size(); };

template<class T>
concept HasSizeAndCopy =
    HasSize<T> && std::is_copy_constructible_v<T>;

template<class T> requires HasSize<T>
int classify(const T&) { return 1; }

template<class T> requires HasSizeAndCopy<T>
int classify(const T&) { return 2; }
\end{lstlisting}

Druga koncepcja \emph{wykorzystuje tę samą} koncepcję
\texttt{HasSize}, dlatego kompilator rozpoznaje wspólny
składnik. Jeśli powtórzymy ręcznie nawet identycznie wyglądające
wymaganie opisujące metodę \texttt{size()},
nie mamy automatycznej gwarancji identyczności ograniczeń
atomowych.

\paragraph{Rekomendacja.}
Twórz proste nazwane koncepcje i buduj bardziej szczegółowe
z ich wykorzystaniem. Ułatwia to konserwację kodu i
przewidywanie wyboru przeciążeń.

\section{Koncepcje dla szablonów klas i ich specjalizacji}

Ograniczać można nie tylko funkcje, lecz również szablony klas.

\begin{lstlisting}
#include <concepts>

template<std::integral T>
class Counter {
    T value_{};
public:
    void add(T delta) { value_ += delta; }
    T value() const { return value_; }
};

// Counter<int> a;       // OK
// Counter<double> b;    // ERROR: constraint not satisfied
\end{lstlisting}

Ograniczenia mogą też sterować wyborem częściowej specjalizacji.

\begin{lstlisting}
#include <concepts>
#include <string_view>

template<class T>
struct TypeLabel {
    static constexpr std::string_view value = "other";
};

template<std::integral T>
struct TypeLabel<T> {
    static constexpr std::string_view value = "integer";
};

template<std::floating_point T>
struct TypeLabel<T> {
    static constexpr std::string_view value = "floating";
};

static_assert(TypeLabel<int>::value == "integer");
static_assert(TypeLabel<double>::value == "floating");
\end{lstlisting}

Nie używaj jednak złożonego systemu specjalizacji, jeżeli
prosta funkcja przeciążona lub \texttt{if constexpr}
lepiej wyraża zamysł.

\section{Koncepcje i pakiety parametrów: kontynuacja rozdziału 4}
\label{sec:concepts-variadic}

Koncepcje bardzo dobrze współpracują z \emph{variadic templates}.
Możemy nakładać warunek na każdy element pakietu typów
za pomocą \emph{fold expression}.

\begin{lstlisting}
#include <concepts>

template<class... Ts>
requires (std::integral<Ts> && ...)
constexpr auto sum_integers(Ts... xs) {
    return (0 + ... + xs);
}

static_assert(sum_integers() == 0);
static_assert(sum_integers(1, 2, 3) == 6);
// sum_integers(1, 2.5);    // ERROR: double not integral
\end{lstlisting}

Zapis \texttt{(P<Ts> \&\& ...)} reprezentuje
koniunkcję wymagań dla wszystkich elementów pakietu.
Dla pustego pakietu taki fold logiczny daje \texttt{true};
nie oznacza to jednak, że każde ciało funkcji poradzi sobie
z zerową liczbą argumentów. W przykładzie użyliśmy
\texttt{0} jako elementu neutralnego sumowania.

\subsection{Ograniczony pakiet typów i uniwersalne przekazywanie}

\begin{lstlisting}
#include <concepts>
#include <utility>

template<std::integral... Ts>
void accept_integers(Ts&&... xs) {
    ((void)std::forward<Ts>(xs), ...); // illustrative use
}
\end{lstlisting}

Zapis \texttt{std::integral... Ts} również ogranicza wszystkie
elementy pakietu. Uwaga: przy dedukcji z przekazywanymi referencjami
\texttt{Ts\&\&}, typ \texttt{Ts} może być \emph{referencją},
gdy argument jest lvalue. W takim przypadku
\texttt{std::integral<Ts>} może być fałszywe.
Jeżeli chcemy akceptować obie kategorie wartości, często
właściwsze jest użycie \texttt{std::remove\_cvref\_t}:

\begin{lstlisting}
#include <concepts>
#include <type_traits>
#include <utility>

template<class... Ts>
requires (std::integral<std::remove_cvref_t<Ts>> && ...)
void accept_any_int(Ts&&... xs) {
    ((void)xs, ...); // demonstration; no ownership transfer
}
\end{lstlisting}

Ta różnica jest szczególnie ważna przy projektowaniu generycznych
wrapperów, które mają przyjmować także obiekty \texttt{const}
i referencje.

\section{Concepts a \texttt{if constexpr}, \texttt{static\_assert} i SFINAE}
\label{sec:concepts-porownanie}

Są to narzędzia pokrewne, ale mają różne cele.

\subsection{\texttt{if constexpr}: wybór implementacji}

\begin{lstlisting}
#include <concepts>
#include <string>
#include <string_view>

template<class T>
std::string describe(const T& value) {
    if constexpr (std::integral<T>) {
        return "an integer";
    } else if constexpr (std::floating_point<T>) {
        return "a floating-point value";
    } else {
        return "another type";
    }
}
\end{lstlisting}

\texttt{if constexpr} umożliwia wybranie gałęzi kodu
\emph{wewnątrz} szablonu. Nie tworzy samo przez się
ograniczonego interfejsu. Z kolei \texttt{requires}
może odrzucić niewłaściwe wywołanie jeszcze na etapie
dobierania przeciążeń.

\subsection{\texttt{static\_assert}: własny komunikat}

\begin{lstlisting}
#include <concepts>

template<class T>
void save_number(const T& value) {
    static_assert(std::integral<T>,
                  "save_number expects an integral type");
    // ...
}
\end{lstlisting}

\texttt{static\_assert} jest użyteczne, gdy potrzebujemy
jednoznacznego komunikatu albo sprawdzamy niezmiennik wewnątrz
implementacji. Nie zastępuje jednak ograniczenia uczestniczącego
w rozstrzyganiu przeciążeń.

\subsection{Przed C++20: \texttt{std::enable\_if} i idiom wykrywania}

\begin{lstlisting}
#include <type_traits>

// C++11/14/17 technique:
template<class T,
         std::enable_if_t<std::is_integral_v<T>, int> = 0>
T old_square(T x) {
    return x * x;
}

// C++20:
#include <concepts>
template<std::integral T>
T new_square(T x) {
    return x * x;
}
\end{lstlisting}

Dla nowego kodu C++20 najczęściej korzystniej wybrać
koncepcję: intencja staje się czytelna i spójna z mechanizmem
porządkowania przeciążeń. \texttt{std::enable\_if}
nadal ma znaczenie w bibliotekach wspierających starsze
wersje standardu.

\section{Praktyczny przykład: średnia z zakresu liczb}
\label{sec:concepts-praktyka}

Zamiast pisać funkcję tylko dla \texttt{std::vector<double>},
możemy określić dokładne wymagania wobec \emph{zakresu}.

\begin{lstlisting}
#include <concepts>
#include <cstddef>
#include <optional>
#include <ranges>
#include <vector>
#include <array>

template<class R>
concept NumericInputRange =
    std::ranges::input_range<R> &&
    std::convertible_to<std::ranges::range_reference_t<R>, double>;

template<NumericInputRange R>
std::optional<double> average(R&& range) {
    double sum = 0.0;
    std::size_t count = 0;
    for (auto&& item : range) {
        sum += static_cast<double>(item);
        ++count;
    }
    if (count == 0) return std::nullopt;
    return sum / static_cast<double>(count);
}

int main() {
    std::vector<int> v{2, 4, 6};
    std::array<double, 2> a{3.0, 7.0};
    auto x = average(v); // 4.0
    auto y = average(a); // 5.0
    auto z = average(std::vector<int>{}); // empty optional
    return !(x && y && !z && *x == 4.0 && *y == 5.0);
}
\end{lstlisting}

Funkcja działa dla różnych kontenerów i zakresów dostarczających
iterację jednokierunkową. Warunek \texttt{input\_range}
wyraża strukturę danych, a \texttt{convertible\_to}
-- oczekiwany rodzaj elementów. Zastosowaliśmy
\texttt{std::optional<double>}, aby pusty zakres nie powodował
dzielenia przez zero. Dla zastosowań numerycznych należy
dodatkowo rozważyć stabilność sumowania, przepełnienia i
znaczenie konwersji do \texttt{double}.

\paragraph{Kiedy \texttt{ranges}?}
Koncepcje standardowych zakresów są szczególnie użyteczne,
gdy funkcja operuje na sekwencji, ale nie wymaga konkretnie
\texttt{std::vector} czy \texttt{std::list}. Zamiast wymuszać
silniejsze wymaganie \texttt{random\_access\_range},
wybieraj najsłabszą koncepcję, która wystarcza do działania
algorytmu.

\section{Koncepcje dla lambd i algorytmów}

Funkcje anonimowe również mogą mieć ograniczone parametry.

\begin{lstlisting}
#include <concepts>

int main() {
    auto twice = [](std::integral auto x) { return x + x; };
    return twice(5) != 10;
}
\end{lstlisting}

Można też tworzyć generyczne algorytmy przyjmujące funkcje
spełniające \texttt{std::predicate}.

\begin{lstlisting}
#include <concepts>
#include <cstddef>
#include <vector>

template<class T, class Pred>
requires std::predicate<Pred&, const T&>
std::size_t count_if_custom(const std::vector<T>& values,
                            Pred pred) {
    std::size_t count = 0;
    for (const T& x : values)
        if (pred(x)) ++count;
    return count;
}
\end{lstlisting}

Zwróć uwagę na argumenty typu w
\texttt{std::predicate<Pred\&, const T\&>}:
sprawdzają możliwość wywołania \emph{konkretnego rodzaju}
predykatu z \emph{konkretnym rodzajem} argumentu.
Kategoria wartości ma znaczenie i nie należy bezrefleksyjnie
wstawiać tu samych nazw typów.

\section{Granice gwarancji: najczęstsze pułapki}
\label{sec:concepts-pulapki}

\subsection{Concept nie sprawdza dowolnej własności semantycznej}

\begin{lstlisting}
#include <concepts>

struct Strange {
    int value{};
    bool operator==(const Strange&) const { return false; }
};

static_assert(std::equality_comparable<Strange>);
// Syntactically valid, but equality lacks reflexivity.
\end{lstlisting}

Kompilator akceptuje składnię porównania, lecz
\texttt{a == a} jest zawsze fałszywe. Kod nie spełnia
matematycznych oczekiwań związanych z relacją równości.
To obowiązek autora typu, niezależny od samego mechanizmu
\texttt{concepts}.

\subsection{Koncepcje nie zamieniają błędów w poprawny kod}

Wyrażenia \texttt{requires} pozwalają w kontekście zależnym
od parametru szablonu traktować niepoprawne podstawienia
jako niespełnione wymagania. Nie oznacza to jednak, że
\emph{każdy} błędny fragment C++ staje się automatycznie
\texttt{false}. Niezależne od parametru, trwale niepoprawne
konstrukcje mogą prowadzić do błędu kompilacji.

\subsection{Różnica między \texttt{T}, \texttt{T\&} i \texttt{const T\&}}

\begin{lstlisting}
#include <concepts>
#include <type_traits>

template<class T>
concept PlainInteger =
    std::integral<std::remove_cvref_t<T>>;

static_assert(!std::integral<int&>);
static_assert(PlainInteger<int&>);
static_assert(PlainInteger<const int&>);
\end{lstlisting}

Przy dedukcji typu generycznego parametru
\texttt{T\&\&} ta różnica często decyduje o tym,
czy z pozoru poprawne wywołanie spełni koncepcję.

\subsection{\texttt{std::same\_as} czy \texttt{std::convertible\_to}?}

\begin{lstlisting}
#include <concepts>

template<class T>
concept StrictAdd = requires(T a, T b) {
    { a + b } -> std::same_as<T>;
};

template<class T>
concept FlexibleAdd = requires(T a, T b) {
    { a + b } -> std::convertible_to<T>;
};

static_assert(!StrictAdd<short>);  // short + short -> int
static_assert(FlexibleAdd<short>); // int -> short allowed
\end{lstlisting}

Koncepcja \texttt{FlexibleAdd} nie oznacza \emph{bezstratnej}
konwersji; \texttt{int} można niejawnie przekształcić do
\texttt{short}, tracąc informację. Nazwa koncepcji powinna
odzwierciedlać rzeczywistą, a nie postulowaną gwarancję.

\subsection{Funkcja ograniczona nie zawsze jest optymalniejsza}

Koncepcje opisują poprawność i wybór instancji szablonów.
Nie wymuszają \emph{inline}, wektoryzacji, eliminacji kopii
ani lepszego kodu maszynowego. Można je stosować bez obietnicy
przyspieszenia wykonywania programu.

\section{Zasady projektowania: kiedy używać, a kiedy nie?}

\begin{center}
\small
\begin{tabular}{|p{0.34\textwidth}|p{0.55\textwidth}|}
\hline
\textbf{Sytuacja} & \textbf{Zalecane rozwiązanie} \\
\hline
Funkcja generyczna wymaga konkretnej kategorii typów & Nazwany \texttt{concept}, np. \texttt{std::integral} \\
\hline
Wymagania dotyczą istnienia operacji lub typu składowego & Własny \texttt{requires-expression} \\
\hline
Ważne są różne wersje algorytmu dla różnych typów & Ograniczone przeciążenia lub specjalizacje \\
\hline
Jedna implementacja ma kilka gałęzi zależnych od typu & \texttt{if constexpr} \\
\hline
Sprawdzasz założenie wewnątrz ciała funkcji & \texttt{static\_assert} z komunikatem \\
\hline
Weryfikujesz jednakowy warunek dla całego pakietu & Fold \texttt{(Concept<Ts> \&\& ...)} \\
\hline
Potrzebujesz iteracji po sekwencji & \texttt{std::ranges::input\_range} lub silniejszy koncept \\
\hline
Funkcja działa dla jednego, konkretnego typu & Zwykła funkcja często jest prostsza niż szablon \\
\hline
Projekt musi działać także w C++17 & \texttt{type\_traits} i (w razie potrzeby) SFINAE \\
\hline
Wymagasz warunku wykonywania programu, np. \texttt{n>0} & Zwykłe sprawdzenie w czasie działania; koncept go nie zastępuje \\
\hline
\end{tabular}
\end{center}

\section{Podsumowanie i pytania kontrolne}

\begin{enumerate}
\item \texttt{concept} nazywa predykat nad parametrami szablonu,
sprawdzany w czasie kompilacji.
\item \texttt{requires-clause} ogranicza deklarację;
\texttt{requires-expression} jest wyrażeniem badającym wymagania.
\item Istnieją cztery formy wymagań: proste, typu, złożone
i zagnieżdżone.
\item \texttt{std::same\_as}, \texttt{std::convertible\_to},
\texttt{std::integral}, \texttt{std::predicate} i koncepcje zakresów
powinny być zwykle pierwszym wyborem.
\item Ograniczenia mogą wpływać na wybór przeciążenia, ale mechanizm
porządkowania nie jest dowolnym dowodzeniem implikacji logicznej.
\item Nazwane koncepcje warto wykorzystywać ponownie przy
budowaniu koncepcji bardziej szczegółowych.
\item \texttt{std::remove\_cvref\_t<T>} ma często znaczenie
przy koncepcjach typów dedukowanych z referencji przekazujących.
\item Koncepcje opisują przede wszystkim wymagania składniowe,
a dodatkowe prawa semantyczne muszą respektować autorzy typów.
\end{enumerate}

\paragraph{Pytania do samodzielnego sprawdzenia.}
Zastanów się, co kompilator zaakceptuje i dlaczego:

\begin{lstlisting}
#include <concepts>
#include <type_traits>

static_assert(std::integral<bool>);
static_assert(!std::integral<int&>);
static_assert(std::integral<std::remove_cvref_t<const int&>>);
static_assert(!std::same_as<int, long>);

// 1. Why does StrictAdd<short> fail?
// 2. Why can a concept be true for a logically broken ==?
// 3. How do requires-clause and requires-expression differ?
// 4. Which overload of category(5) is chosen?
// 5. Why is NumericInputRange more general than vector<double>?
\end{lstlisting}

\paragraph{Źródła i dalsza lektura.}
Oficjalne reguły języka określa standard ISO C++, a przystępne
opisy i przykłady znajdują się w serwisie \emph{cppreference}.
Przydatne adresy:
\begin{lstlisting}
en.cppreference.com/cpp/language/constraints
en.cppreference.com/cpp/language/requires
en.cppreference.com/cpp/concepts
\end{lstlisting}
